\documentclass[11pt]{article}
\usepackage[margin=1in]{geometry}
\usepackage{amsmath,amssymb,amsfonts,amsthm,mathtools,bm}
\usepackage[T1]{fontenc}
\usepackage{microtype}
\newtheorem{theorem}{Theorem}
\newtheorem{lemma}{Lemma}
\newtheorem{cor}{Corollary}
\newtheorem{defn}{Definition}
\newtheorem{principle}{Principle}
\newtheorem{statement}{Statement}
\theoremstyle{remark}
\newtheorem{statementr}{Statement}

\title{Identical Particles, Set-Indexed Spin Bundles, and Bundle Triviality}
\author{James Taylor}
\date{Working draft derived from the 2003 thesis}
\begin{document}
\maketitle

\section{Many particle systems} \label{c:manpart}


This chapter  
demonstrates that Bohmian mechanics\  provides an argument for why
the wave function\ for a system of identical particles must have a certain symmetry;
this should be contrasted with the claim of some that the lack of precise
trajectories in standard quantum mechanics is the reason for such a
constraint. 
 There are others who have discussed the Bohmian\ story of
identical particles \cite{SC95, S95, BBS99, Bac03}. Our
particular approach has the appeal that it is 
founded upon the Abelian Quantization Principle. 


The principle implies multiple possibilities only for 
configuration spaces with nontrivial topology.
As we shall discuss, for $N$ identical particles, the 
configuration space, \ensuremath{{}^N\ensuremath{\mathbb{R}^{3}}},  is the space of all $N$-element subsets of
physical space. The space is a manifold which is not simply connected;
its fundamental group has two characters. Applying the principle leads
to 
the Bose-Fermi alternative.
In order to
apply our principle in 
the case of spin, we define 
the spin bundle over \ensuremath{{}^N\ensuremath{\mathbb{R}^{3}}}.  It is an
unusual definition;  the fiber at the point $\ensuremath{q} \subset
\ensuremath{\mathbb{R}^{3}}$, is the $N$-fold tensor product of the 1-particle spin
space with itself using
\ensuremath{q}\ as the index set for 
the
tensor product. Once we have the spin bundle, then the principle
immediately leads to the Bose-Fermi alternative. 


The literature on identical particles is quite extensive.\footnote{We found one
compilation of references for the symmetrization postulate that
contained over 350 references and yet it seemed far from complete.}
 We will only mention some of the key papers that we managed to come
across.
The literature about the symmetrization postulate
in non-relativistic quantum mechanics seems to have begun with
\cite{MO64}. Both that paper and \cite{Gir65} focussed on 
observables. This approach is still used; in  \cite{Gal91}, there
is an argument for the identical particle case 
based on observables and decomposition of the Hilbert space into
permutation representations. A critique of such approaches may be
found in \cite{Mir73}. 
The topological approaches began with \cite{DL71} and its use of Feynman path
integrals on 
multiply
connected spaces. 
An approach that changes the connection on the configuration space is
that of \cite{LM77}; they include a brief discussion for the case of spin. 
An orthodox approach, using the covering space and only discussing
scalar wave functions,  
is discussed in \cite{Mor92}.  
The version for
Nelson's stochastic 
mechanics may be found in Nelson 
\cite{Nel85}

and its relevance is that it is that description upon which this
chapter is based. 


We do not address the connection
between spin and statistics. Our accomplishment is to demonstrate that
for each type of particle, there are two choices for theories that
locally look identical; the spin-statistics is an additional statement
that gives the relation between spin and statistics. It is usually
derived in the context of relativistic quantum field theory.
In \cite{Bal93}, the authors give
an approach that does establish the relationship between spin and
statistics in the context of non-relativistic theories describing the
creation and annihilation of particles with 
anti-particles. The approach is based
on making a radical change to the 
configuration space. Indeed, it ceases to be a manifold and they 
devote many pages to arguing that the topology is reasonable and
simply connected. 
For scalar wave functions, this implies that only bosonic wave functions are
allowed. For spin, they specify how two particles can annihilate or
create one another; this mimics a full rotation in physical space
acting on a 1-particle spin space thereby providing 
the spin-statistics relation.
In \cite{MR97}, the authors use the Schwinger representation of spin to
give an argument for the spin-statistics relation in non-relativistic
quantum mechanics. They postulated that a certain set of axioms, which 
their
specific construction did satisfy, would force the spin-statistics
relation to be true. The conjecture is false as  they 
later found other constructions satisfying these axioms, but 
 not the
spin-statistics relation \cite{MR01}. 


This chapter is organized as follows. We start by explaining \ensuremath{{}^N\ensuremath{\mathbb{R}^{d}}}. We
follow this by describing a set of cuts for that space. We then apply
the Abelian Quantization Principle\ to scalar-valued wave functions to  
conclude that
the particles are either governed by bosonic or fermionic wave functions if
$d \geq 3$. Anyons appear as possibilities if $d = 2$.  
We make full use of the notion of a set-indexed tensor product  
in Section {\ref{s:nspin}}\  when we discuss spin
on the natural 
configuration space. Our principle immediately applies and we obtain
the usual Bose-Fermi alternative in the case of spin. Section {\ref{s:bprin}}\
considers, and then rules out, certain possibilities that go beyond the
general principle. In the specific case of identical particles, we
give, in Section {\ref{s:view}}, a different 
 kind of argument for the Bose-Fermi
alternative. Section {\ref{s:cm}}\  briefly  highlights why classical mechanics,
involving identical particles and formulated on \ensuremath{{}^N\ensuremath{\mathbb{R}^{d}}},
are not enriched by appealing to the covering space if $d \geq 3$. 
The final section describes a speculative idea as to how to describe a
system of distinguished particles as arising from a theory of
indistinguishable particles.


\section{The configuration space of identical particles} \label{s:nat}


This section explains why
the space of $N$-element subsets of physical space is the correct
configuration space. The usual way of dealing with identical particles
is to label them and then demand that the labelling is irrelevant to
the motion of the particles. The configuration space which we use, and
which is certainly not novel, is a way of formulating the dynamics
without introducing the irrelevant structure of labelling.  In what
follows, we shall take $\ensuremath{\mathbb{R}^{\ensuremath{d}}}$ to be the physical space and we
require $\ensuremath{d} \geq 2$.\footnote{For $\ensuremath{d} =1$, the configuration space
consists of a variety of simply connected components. The principle
will not apply and the statistics, in as much as they are present, 
 only arise as boundary conditions on the Hamiltonian.}


As a starting point,  we might say that a 
 dynamical system of $N$ identical particles is one in which the
appropriate exchange
of the initial 
conditions results in the new solutions being the exchange of the old
solutions.  
Note that to exchange the particles, we have already
distinguished them by more than just their locations. The statement is
then asserting that the dynamics does not notice the distinguishing
characteristics. A stronger statement, and one more to the point, is
 that
exchanging two identical particles has no effect on the world. 
It would then be logical to make no distinction between the particles
in forming the configuration space of the system. 
Fundamentally, we have $N$ points in
physical space. 
Stating this mathematically, our configuration space is 
the space of all $N$-element subsets of \ensuremath{\mathbb{R}^{\ensuremath{d}}}.  Instead of 
\ensuremath{(\boldsymbol{q}_1,\ldots,\boldsymbol{q}_N)}, the configuration is  
\ensuremath{\{\boldsymbol{q}_1,\ldots,\boldsymbol{q}_N\}}. Instead of \ensuremath{\mathbb{R}^{\ensuremath{d} N}}, the  configuration space of $N$
identical particles is given by 
\begin{equation*}
\ensuremath{{}^N\ensuremath{\mathbb{R}^{\ensuremath{d}}}} \  \ensuremath{:=} \ \ensuremath{\{S|S \subseteq \ensuremath{\mathbb{R}^{\ensuremath{d}}}, |S| = N \}}. 
\end{equation*}

A more familiar form is given by the following; it is the mathematical
version of starting with a labelled system and then removing the labelling.  
We define $\ensuremath{\ensuremath{\Delta}}$ to be the set of configurations of \ensuremath{\mathbb{R}^{\ensuremath{d} N}}\ having
coincident points, that is to say,  
\begin{equation*}
\ensuremath{\ensuremath{\Delta}} \ensuremath{:=} \ensuremath{\{ \ensuremath{(\boldsymbol{q}_1,\ldots,\boldsymbol{q}_N)} \in \ensuremath{\mathbb{R}^{\ensuremath{d} N}} | \exists \text{
$i$, $j$ s.t. $\boldsymbol{q}_i = \boldsymbol{q}_j$} \}}. 
\end{equation*}  
let $\ensuremath{\mathbb{R}^{\ensuremath{d},N}}_{\neq}$
 denote the set of noncoincident points in \ensuremath{\mathbb{R}^{\ensuremath{d} N}}, i.e.\ 
\begin{equation*}
\ensuremath{\mathbb{R}^{\ensuremath{d},N}}_{\neq} \ensuremath{:=}\ \ensuremath{\ensuremath{\mathbb{R}^{\ensuremath{d} N}} \setminus \ensuremath{\ensuremath{\Delta}}}. 
\end{equation*}
Let \ensuremath{S_N}\ denote the group of permutations of the set 
$\{ 1, \ldots N \}$. 
This group acts freely on \ensuremath{\ensuremath{\mathbb{R}^{d,N}}_{\neq}}\ by permuting the indices, i.e.\ 
\begin{equation*}
\ensuremath{\sigma} (\boldsymbol{q}_1,\ldots, \boldsymbol{q}_N) \mapsto 
(\boldsymbol{q}_{\ensuremath{\sigma}1},\ldots, \boldsymbol{q}_{\ensuremath{\sigma} N})
\end{equation*} 
and the exclusion of \ensuremath{\ensuremath{\Delta}}\ implies that, for every nontrivial
  permutation, 
there are no fixed points.
Then the natural configuration space \ensuremath{{}^N\ensuremath{\mathbb{R}^{\ensuremath{d}}}}\ can be (canonically)
identified with the Riemannian manifold of equivalence classes under
that action; more compactly
\begin{equation*}
 \ensuremath{{}^N\ensuremath{\mathbb{R}^{\ensuremath{d}}}} \ensuremath{\cong} \ensuremath{\ensuremath{\mathbb{R}^{d,N}}_{\neq}} / \ensuremath{S_N}.
\end{equation*}
From this identification, \ensuremath{{}^N\ensuremath{\mathbb{R}^{\ensuremath{d}}}}\ is itself a Riemannian manifold. 
The identification also correctly suggests that, in analogy with the
torus, \ensuremath{{}^N\ensuremath{\mathbb{R}^{\ensuremath{d}}}}\ 
is not simply connected. 


\section{Applying the Abelian Quantization Principle}


For $\ensuremath{d} \geq 3$, the universal covering
space of \ensuremath{{}^N\ensuremath{\mathbb{R}^{\ensuremath{d}}}}\ is 
 \ensuremath{\ensuremath{\mathbb{R}^{d,N}}_{\neq}}. As said above, a permutation naturally defines a
diffeomorphism on \ensuremath{\ensuremath{\mathbb{R}^{d,N}}_{\neq}}; the 
set of these diffeomorphisms forms the covering group. 
The fundamental group at a point $\ensuremath{q}\ \in \ensuremath{{}^N\ensuremath{\mathbb{R}^{\ensuremath{d}}}}$ is the set of
bijections of the set to itself. Both groups are isomorphic to the
permutation group \ensuremath{S_N}. 
The only characters of \ensuremath{S_N}\  are the trivial one and
the alternating 
one. The trivial character defines the bosonic wave functions which satisfy
the periodicity condition
\begin{equation} \label{e:tri}
\ensuremath{\psi}(\ensuremath{\sigma} \ensuremath{\hat{q}}) = \ensuremath{\psi}(\ensuremath{\hat{q}}).
\end{equation}
The alternating character defines the fermionic wave functions which satisfy
the periodicity condition 
\begin{equation} \label{e:alt}
\ensuremath{\psi}(\ensuremath{\sigma} \ensuremath{\hat{q}}) = (-1)^{\ensuremath{\sigma}}\ensuremath{\psi}(\ensuremath{\hat{q}}).\footnote{
The sign of the permutation, $(-1)^\ensuremath{\sigma}$, is  $1$ if
\ensuremath{\sigma}\ can be written as a composition of an even number of
transpositions, and $-1$ otherwise.}
\end{equation} 
Applying the Abelian Quantization Principle, we conclude the Bose-Fermi alternative.

The case  $d = 2$ is quite different since $\ensuremath{\mathbb{R}^{2N}}_{\neq}$ is not
simply connected. The fundamental group of
 \ensuremath{{}^N\ensuremath{\mathbb{R}^{\ensuremath{2}}}}\ is isomorphic to the braid group, not
\ensuremath{S_N}. There is a 
one-parameter family of different characters for this group. 
Bosons and fermions are among the possibilities, but there are many
more. The reason that $d = 2$ is different 
is that a set of paths that  exchanges two
particles twice may not be homotopic to the identity; more to the
point, these paths are in $\ensuremath{\mathbb{R}^{2}}\times[0,1]$, 
a space in which knots can
occur. 
The generators for
the braid group may be chosen to be a certain subset of 
 braids that exchange two particles;\footnote{More precisely, label
the particles. Then $\sigma_i$ is the path exchanging particles
$i$ and $i+1$. Furthermore, these paths, which  may be thought of as
parametrized 
circles in the plane, should all have the same orientations arising
from the parametrization.} 
they satisfy the relations
\begin{align*} 
\sigma_i \sigma_j &=\sigma_j\sigma_i \quad \hbox{for} \quad j \neq i \pm 1 , \\
\sigma_i\sigma_{i+1}\sigma_i &= \sigma_{i+1}\sigma_i \sigma_{i+1}. 
\end{align*}
The second relation implies that a character of the braid group assigns the same complex number to
each  generator. If that number is $e^{i \beta}$, 
then the periodicity condition is 
\begin{equation*}
\ensuremath{\psi} (\ensuremath{\sigma}(\ensuremath{\hat{q}}))
= e^{i |\ensuremath{\sigma}| \beta}  \ensuremath{\psi}(\ensuremath{\hat{q}}),
\end{equation*} 
where $|\ensuremath{\sigma}|$ is the, appropriately counted, number of generators
defining  $\ensuremath{\sigma}$. Bosons correspond to $\beta = 0$ and fermions to
$\beta = \pi$. 
The particles corresponding to the  other possibilities are
usually called \emph{anyons}. They were first suggested in
\cite{LM77}, and their investigation began in earnest with 
\cite{GMS81, Wi82}. See \cite{Mor92} for some more details and references. 

\subsection{Identical particles moving in $M$}

More generally, let physical space be the Riemannian manifold $M$. Then
the configuration space for identical particles would be the
$N$-element subsets of $M$, denoted by \ensuremath{{}^{N}\! {M}}. The 
tangent space of
the configuration space at the point \ensuremath{q}\ 
is $\bigoplus_{\ensuremath{\boldsymbol{q}} \in \ensuremath{q}} T_{\ensuremath{\boldsymbol{q}}}M$. Its metric, if the
particles all have mass $m$, is
\begin{equation*} 
\ensuremath{{}^{N}\! {g}}(\mathop{\oplus}_{\ensuremath{\boldsymbol{q}} \in \ensuremath{q}} v_{\ensuremath{\boldsymbol{q}}}, \mathop{\oplus}_{\ensuremath{\boldsymbol{q}}
\in \ensuremath{q}} w_{\ensuremath{\boldsymbol{q}}})
\ensuremath{:=} \sum_{\ensuremath{\boldsymbol{q}} \in \ensuremath{q}} m g(v_{\ensuremath{\boldsymbol{q}}},w_{\ensuremath{\boldsymbol{q}}}).
\end{equation*} 

Then, as above, we can apply the principle to \ensuremath{{}^{N}\! {M}}. 
We will always have bosons and fermions as
possibilities, but there may be more possibilities depending on the
dimension of $M$ and $M$'s fundamental group.  


\section{Identical particles with spin} \label{s:nspin}


The textbooks, when discussing identical particles with spin $s$, take 
the configuration space to be \ensuremath{\ensuremath{\mathbb{R}^{d,N}}_{\neq}}\ and the value space of the wave function\ to
be the $N$-fold tensor product of $W \ensuremath{:=} \ensuremath{\mathbb{C}}^{2s+1}$, the
1-particle value space, with itself. We denote it in various ways:
$
 W_1 \otimes \cdots \otimes W_N 
=  \bigotimes_{i=1}^{N} W_i = \bigotimes_{i \in 
\{1, \ldots, N \}} W_i = W^{\otimes N}.
$

It is tempting to say that each factor of the tensor product
corresponds to each particle. But if we do this, we are
assigning a label  $i \in \{1, \ldots, N\}$ to each particle. On \ensuremath{{}^N\ensuremath{\mathbb{R}^{3}}},
there is no consistent way of making such an assignment. 
Thus it may appear that the configuration space \ensuremath{{}^N\ensuremath{\mathbb{R}^{3}}}\ is incompatible
with spin, but it is, in fact, quite compatible. 
We associate a copy
of $W$ not with each particle, but rather with every
point in physical space $\ensuremath{\mathbb{R}^{3}}$. 
Even in the
standard story,  labelling identical particles is not allowed. 
That is why, even for bosons, there is a permutation
action on the value space. If we wish to work on
\ensuremath{{}^N\ensuremath{\mathbb{R}^{3}}}, we must define a tensor product which only labels the particles
by their locations.
We do this by replacing the usual  tensor product's 
index set $\{1, \ldots, N\}$  
with the $N$-element set $\{ \ensuremath{\boldsymbol{q}}_1, \ldots \ensuremath{\boldsymbol{q}}_N\} \in \ensuremath{{}^N\ensuremath{\mathbb{R}^{3}}}$. 
Instead of using
$\bigotimes_{i \in 
\{1, \ldots, N \}} W_i = W^{\otimes N},
\text{ we form } \bigotimes_{\ensuremath{\boldsymbol{q}} \in 
\ensuremath{q}} W_{\ensuremath{\boldsymbol{q}}} = W^{\otimes \ensuremath{q}}.
$

\subsection{The spin bundle and a Bohmian theory}
The use of an arbitrary finite index set is a straightforward notion. 
We wish to replace $\{1, \ldots, N\}$ with the $N$-element set
$T$. The ordering in the usual tensor product serves as a way to map
the factors of the tensor product bijectively to the elements of $\{1,
\ldots, N\}$.  Analogously, a product element in the  $T$-indexed tensor
product of $W$ with itself $N$-times 
should be thought of as a function from $T$ into $W$. The obvious
analogues to the usual definitions of the tensor product hold and we
state them in Appendix \ref{s:setten}; one can also find this
definition in  \cite{Bour}. Here we shall content ourselves to make a
notational correspondence. The tensor product space $W^{\otimes N 
}$ is replaced with $W^{\otimes T}$. The product element
$\otimes_{i=1}^N w_i$ is replaced with $\otimes_{a \in T} w_a$. The usual
tensor product rules still hold. For example, 
$ \mathop{\otimes}_{a \in T} x_a + \mathop{\otimes}_{a \in T} y_a = \mathop{\otimes}_{a
\in T} z_a$ 
if, for some $a' \in T$, we have both $x_{a'}
+y_{a'} = z_{a'}$ and, for all the other $a \in T$, 
 $x_a = y_a=z_a$.
This trivial extension of the commonly used version of the tensor 
product is exactly what we need for spin. 
 


We begin by defining the spin bundle,\footnote{For spin, we shall
 restrict our discussion to $d=3$. We will use the
 notation

$\ensuremath{\mathcal{Q}} \ensuremath{:=} \ensuremath{{}^N\ensuremath{\mathbb{R}^{\ensuremath{3}}}} $

and use \ensuremath{\widehat{\ensuremath{\mathcal{Q}}}}\ for the universal covering space, \ensuremath{\ensuremath{\mathbb{R}^{d,N}}_{\neq}}.
  }
which we shall write as
\ensuremath{\ensuremath{W}^{\otimes \ensuremath{\mathcal{Q}}}}. Let \ensuremath{W}\ be the single-particle 
spin space for  
our $N$ identical particles.
 Given $\ensuremath{q}
\in \ensuremath{\mathcal{Q}}$, we define the spin space \emph{at that point} to be
\ensuremath{{\ensuremath{W}}^{{\otimes}{\ensuremath{q}}}}. This defines the fibers for our spin bundle; one can use the
connection described below to define the bundle
structure.\footnote{One could always define the bundle structure first
by locally labeling the particles in order to define a coordinate
  patch, but it is more appealing to use the
identification between the points that the connection provides.}

We need to define the connection. We shall do this by  
specifying the equivalent notion of parallel transport. Given
two points \ensuremath{q}, \ensuremath{q'}\ and a path \ensuremath{\alpha}\ connecting them, then
we have a  set bijection \ensuremath{\ensuremath{\sigma}_{\alpha}}\ between \ensuremath{q}\ and
\ensuremath{q'}\ given by following the $N$ trajectories in physical space that
\ensuremath{\alpha}\ represents. This bijection of the index sets provides a
natural  isomorphism between \ensuremath{{\ensuremath{W}}^{{\otimes}{\ensuremath{q}}}}\ 
and \ensuremath{{\ensuremath{W}}^{{\otimes}{\ensuremath{q}'}}}; this isomorphism will be the parallel transport operator,
or holonomy operator, 
associated to \ensuremath{\alpha}. It is clear that it depends  only on the
homotopy class of \ensuremath{\alpha}.
This in turn implies that the connection is flat. We shall use
$\nabla$ to denote the gradient associated with this
connection.


Section {\ref{s:bbun}} has essentially already described the 
Bohmian mechanics on this bundle. Indeed, the wave function\ is a
(smooth) section
of the bundle. This  means that $\ensuremath{\psi}( \ensuremath{q}) \in \ensuremath{{\ensuremath{W}}^{{\otimes}{\ensuremath{q}}}}$ for
every $\ensuremath{q} \in \ensuremath{\mathcal{Q}}$. In particular, the value space of the wave function\
varies from point to point. 
The wave function\ defines a velocity field on
\ensuremath{\mathcal{Q}}\ as it did before. The equations
of motion are Bohm's equation and Pauli's equation:
\begin{subequations}
\begin{align} 
\ensuremath{\frac{d{\boldsymbol{Q}_t}}{d{t}}} =&  \ensuremath{\frac{\hbar}{m_{}}} \textrm{Im}
\frac {\ensuremath{({\ensuremath{\psi}_t}, {\nabla_{\boldsymbol{Q}_t} \ensuremath{\psi}_t})}}{
\ensuremath{({\ensuremath{\psi}_t}, {\ensuremath{\psi}_t})}} (Q_t)
  \quad \text{for each $\boldsymbol{Q}_t \in Q_t$,} \label{e:natbohm}\\  
i \hbar \ensuremath{\frac{\partial {\ensuremath{\psi}_t}}{\partial {t}}} =& (\ensuremath{ - \sum_{ \ensuremath{\boldsymbol{q}} \in \ensuremath{q}}
\ensuremath{\frac{\hbar^2}{2m_{}}} \ensuremath{\Delta}_{\ensuremath{\boldsymbol{q}}}  + V_{\ensuremath{q}} + \sum_{\ensuremath{\boldsymbol{q}} \in \ensuremath{q}}
\ensuremath{\ensuremath{\mu}(\ensuremath{\boldsymbol{S}}^{\ensuremath{\boldsymbol{q}}} \cdot
\boldsymbol{B}_{{\ensuremath{\boldsymbol{q}}}})}}) \ensuremath{\psi}_t \label{e:natschr}.
\end{align}
\end{subequations}
By $\nabla_{\boldsymbol{Q}} \ensuremath{\psi}$, we mean that we fix all the elements of
$Q$ except for $\boldsymbol{Q}$.\footnote{More precisely, $\nabla \ensuremath{\psi}$ is
metrically equivalent to the
1-form $D \ensuremath{\psi}$ acting on 
$T_{\ensuremath{q}}\ensuremath{\mathcal{Q}} = \bigoplus_{\ensuremath{\boldsymbol{q}} \in \ensuremath{q}} T_{\ensuremath{\boldsymbol{q}}}
\ensuremath{\mathbb{R}^{\ensuremath{3}}}$. 
As we want to transform a 1-form defined on $T_{\boldsymbol{Q}}
\ensuremath{\mathbb{R}^{\ensuremath{3}}}$, we restrict $D \ensuremath{\psi}$ to the subspace in
$T_{\ensuremath{q}}\ensuremath{\mathcal{Q}}$ 
that   $T_{\boldsymbol{Q}}
\ensuremath{\mathbb{R}^{\ensuremath{3}}}$ is. Equivalently, we are transforming the 1-form
on $T\ensuremath{\mathcal{Q}}$ defined by
\mbox{$D\ensuremath{\psi} \{ (\bigoplus_{\substack{\boldsymbol{Q}' \in Q, 
\boldsymbol{Q}' \neq  \boldsymbol{Q}}} w_{\boldsymbol{Q}'}) \oplus
v_{\boldsymbol{Q}} \}\ensuremath{:=} D\ensuremath{\psi}  
\{0 \oplus
v_{\boldsymbol{Q}} \}$.}
}
The symbol
$\ensuremath{\frac{d{\boldsymbol{Q}_t}}{d{t}}}$ represents  the tangent vector of the
trajectory in physical 
space going through $\boldsymbol{Q}_t$ at time $t$. 


 


\subsection{Applying \ensuremath{\text{AQP}_{\ensuremath{\widehat{\ensuremath{\mathcal{Q}}}}}}}

The covering space version of the Abelian Quantization Principle\ 
tells us to lift the spin bundle to the covering
space and look at wave functions satisfying a periodicity condition. 
When we pull-back the bundle, by definition
elements in the same covering fiber\ 
have the same exact spin space. We shall  use the notation \ensuremath{\ensuremath{W}^{\otimes \ensuremath{\widehat{\ensuremath{\mathcal{Q}}}}}}\ for
 the lifted spin 
bundle since  the lifted bundle is identical to
 the bundle formed using the index set
 $\{\ensuremath{\boldsymbol{q}}_1, \dotsc \ensuremath{\boldsymbol{q}}_N\} = \ensuremath{q}$ at 
 the point $(\ensuremath{\boldsymbol{q}}_1, \dotsc, \ensuremath{\boldsymbol{q}}_N)= \ensuremath{\hat{q}}$. We shall also use the
 notation $\ensuremath{\boldsymbol{q}} \in \ensuremath{\hat{q}}$ with the idea of viewing $\ensuremath{\hat{q}}$ as a set
 instead of as an $N$-tuple. 

The class of wave functions for each theory generated by the principle 
 is characterized by the periodicity condition 
\begin{equation*}
\ensuremath{\psi} (\ensuremath{\sigma}(\ensuremath{\hat{q}}))
= \ensuremath{\gamma_{\ensuremath{\sigma}}} \ensuremath{\psi}(\ensuremath{\hat{q}}).
\end{equation*} 
 Since
 our space is \ensuremath{{}^N\ensuremath{\mathbb{R}^{\ensuremath{3}}}}, we have two 
 possibilities for different 
 classes  of wave functions: bosons  ($\ensuremath{\gamma_{\ensuremath{\sigma}}} = 1$)  and fermions ($\ensuremath{\gamma_{\ensuremath{\sigma}}}
 =(-1)^{\ensuremath{\sigma}}$).  We have reached the main
 conclusion.  


It is useful to relate our presentation to the usual story. 
In the usual story, one starts with
 \ensuremath{\ensuremath{\mathbb{R}^{d,N}}_{\neq}}\ and one  uses the trivial bundle $W^{\otimes N}$ with its trivial
connection. By a bosonic wave function, one then means a wave function\ satisfying
\begin{equation} 
\ensuremath{\psi}(\ensuremath{\sigma} \ensuremath{\hat{q}}) = \ensuremath{R_{{\ensuremath{\sigma}}}} \ensuremath{\psi}(\ensuremath{\hat{q}}).
\end{equation}
where \ensuremath{R_{{\ensuremath{\sigma}}}}\ is the permutation of the spin factors related to the
permutation of the set $\{ 1, \ldots, N \}$ represented by
$\ensuremath{\sigma}$.
To
relate our spin space, $\ensuremath{{\ensuremath{W}}^{{\otimes}{\ensuremath{\hat{q}}}}}$, to the trivial space,
$\ensuremath{{\ensuremath{W}}^{{\otimes}{N}}}$, at the point $\ensuremath{\hat{q}} = (\ensuremath{\boldsymbol{q}}_1, \dotsc, \ensuremath{\boldsymbol{q}}_N)$, we use
the bijection between $\{1, \dotsc, N\}$ and $\{\ensuremath{\boldsymbol{q}}_1, \dotsc,
\ensuremath{\boldsymbol{q}}_N\}$ which the $N$-tuple provides. 

We would like to point out that our general principle would not apply
in the usual formulation. Indeed, if we used the trivial bundle, and
its trivial connection, on the
base space, we first run into the problem that the Hamiltonian is not
particularly well-defined.\footnote{Note that we
do know the Hamiltonian on the covering space; the spin terms in the
Hamiltonian at the points $\ensuremath{\hat{q}}$ and $\ensuremath{\sigma} \ensuremath{\hat{q}}$,  are related by
conjugation with \ensuremath{R_{{\ensuremath{\sigma}}}}. 
} 
If we then lift the
trivial bundle  $\ensuremath{\mathcal{Q}} \times \ensuremath{{\ensuremath{W}}^{{\otimes}{N}}}$ to \ensuremath{\widehat{\ensuremath{\mathcal{Q}}}}, 
we get the trivial flat bundle,
$\ensuremath{\widehat{\ensuremath{\mathcal{Q}}}}\times \ensuremath{{\ensuremath{W}}^{{\otimes}{N}}}$. As implied
in Section {\ref{s:bprin}}, sections of $\ensuremath{\widehat{\ensuremath{\mathcal{Q}}}} \times \ensuremath{{\ensuremath{W}}^{{\otimes}{N}}}$
  which satisfy the trivial periodicity
condition will not be
invariant under the evolution. 
Only the fermionic and bosonic periodicity conditions are invariant
under the Pauli evolution. 

\subsection{Applying \ensuremath{\text{AQP}_{H}}} 


As for the cut version, we have the same conclusions as before and for
the same reason. For the cuts of \ensuremath{{}^N\ensuremath{\mathbb{R}^{\ensuremath{3}}}}\ described in
Section {\ref{s:cutnrd}}, we have the 
boundary condition of $-1$ for the fermions. The BV-equivalence maps
consist of multiplication by $1$ or $-1$ depending on the cut and the
regions. This approach suggests that the difference
between bosons and fermions is the Hamiltonian. In particular, a wave function,
whose support is simply connected, has the potential to be either
bosonic or fermionic depending on which evolution is
chosen. Also note that, just as an example, the  other fermionic wave functions will
evolve under the bosonic evolution.  
But  they  
 should instantaneously cease to be fermionic as well as becoming
 unsuitable for a Bohmian evolution.

\subsection{Cutting \ensuremath{{}^N\ensuremath{\mathbb{R}^{\ensuremath{3}}}}} \label{s:cutnrd}

A rather pleasant exercise  is to 
describe how to cut for $\ensuremath{\mathcal{Q}} \ensuremath{:=} \ensuremath{{}^N\ensuremath{\mathbb{R}^{\ensuremath{3}}}}$ so that it becomes simply
connected.  
For each $i$, we define $\ensuremath{\kappa}_i$, $i=1,\dotsc,\ensuremath{3}$, to
be the set of points $\ensuremath{q} \in \ensuremath{\mathcal{Q}}$ such that at least two points
in \ensuremath{q}\ have equal  $\ensuremath{{i}^{th}}$\  components. In other words, 
\begin{equation*}
\ensuremath{\kappa}_i \ensuremath{:=} \{ \ensuremath{q} \in \ensuremath{\mathcal{Q}}| \exists \ensuremath{\boldsymbol{q}}, \tilde{\ensuremath{\boldsymbol{q}}} \in \ensuremath{q}
\text{ s.t. } \ensuremath{\boldsymbol{q}}^i =\tilde{\ensuremath{\boldsymbol{q}}}^i\}.
\end{equation*} 
Note that $\ensuremath{q} \in \ensuremath{\bar{\ensuremath{\mathcal{Q}}}_{{\ensuremath{\kappa}_i}}}$ iff \ensuremath{q}\ can be ordered
by the $\ensuremath{{i}^{th}}$\ component of its
elements. 
Indeed, a corresponding fundamental domain in \ensuremath{\widehat{\ensuremath{\mathcal{Q}}}}\ is the set
$\{\ensuremath{\hat{q}} \in \ensuremath{\widehat{\ensuremath{\mathcal{Q}}}}| \ensuremath{\boldsymbol{q}}_1^i < \ensuremath{\boldsymbol{q}}_2^i < \cdots < \ensuremath{\boldsymbol{q}}_{N}^i
\}$. From the definition of a cut, it is reasonably clear that
properties 1 and 2 are satisfied;  the main observation is that
the paths must preserve the ordering. Thus nontrivial exchange is
forbidden. As for the third property, note that the set in which
three points have the same $\ensuremath{{i}^{th}}$ component is a
$(n-2)$-dimensional submanifold. But if we restrict the submanifold to
 a small enough
neighborhood of such a
point, then we find that it becomes a union of 3 regular submanifolds, 
one for each
pair. We could also have a $(n-2)$-dimensional part by having two
pairs of points with the same $\ensuremath{{i}^{th}}$\ component for each pair, but
different for the different pairs, e.g. $\ensuremath{\boldsymbol{q}}_1^i = \ensuremath{\boldsymbol{q}}_2^i \neq
\ensuremath{\boldsymbol{q}}_3^i = \ensuremath{\boldsymbol{q}}_4^i$. Then, in small neighborhoods, this
appears as the union of two submanifolds. If all of the $\ensuremath{{i}^{th}}$\
components are 
equal, then there are $\binom{N}{2}$ submanifolds in the local union. 
In analogous fashion, all of the self-intersections are locally the
union of regular submanifolds. 


\subsection{Applying \ensuremath{\text{AQP}_{\nabla}}}

Instead of viewing fermionic wave functions as sections on the covering space
satisfying a periodicity condition or as sections of the cut spin bundle
satisfying certain virtual boundary conditions, we will view them as
sections of a certain bundle over
\ensuremath{{}^N\ensuremath{\mathbb{R}^{\ensuremath{3}}}}. For
complex-valued functions, we will construct the Fermi line, a one
dimensional bundle endowed with a flat connection such that parallel
transporting a vector along an element \ensuremath{\sigma}\ of the fundamental group 
 is equivalent to
multiplying the vector by $(-1)^{\ensuremath{\sigma}}$. 
As for
higher spins, we construct the appropriate bundle for fermions 
by tensoring the Fermi
line with the spin bundle. 

We start with some terminology. 
 A Bose bundle is a
bundle whose sections are bosonic wave functions; a Fermi bundle's  sections
 are the fermionic wave functions.
For spin 0, i.e.\ the scalar case,  we shall refer to the
Bose bundle as  the Bose
line and the Fermi bundle as the Fermi line. 
 The number of particles for the discussion below will be, as usual, 
 $N$.


The Bose line is the trivial flat line bundle. That is to say the Bose
line is $\ensuremath{\mathcal{Q}} \times \ensuremath{\mathbb{C}}$. 
For higher spins, the Bose bundle is the spin bundle, 
 a flat bundle which is inequivalent to the trivial flat
 bundle. Indeed, the
holonomy representation on the spin bundle is the permutation
 representation.\footnote{   
By  \emph{permutation representation}, we mean any representation
 equivalent to the one acting on $(\ensuremath{\mathbb{C}}^k)^{\otimes
 N}$  which takes any $N$-element permutation into the
 operator \ensuremath{R_{{\ensuremath{\sigma}}}}\ which does the same permutation on the tensor product index
 set $N$. 
The \emph{alternating permutation
representation} shall mean any representation equivalent to $\ensuremath{\sigma}
\mapsto (-1)^{\ensuremath{\sigma}} \ensuremath{R_{{\ensuremath{\sigma}}}}$.
}    
We shall describe the Fermi line in a
moment; assume  that we know what it is. 
The Fermi bundle for higher spins is then the tensor product of the
Bose bundle with the Fermi line as the principle suggests. 
By construction, the holonomy representation of 
the fermion bundle is the alternating permutation representation, as
 it should be.

The question remains how to describe  the Fermi line over \ensuremath{{}^N\ensuremath{\mathbb{R}^{\ensuremath{3}}}}. The
 Fermi line is a flat, 1-dimensional
 bundle such that parallel transport maps   
$w \mapsto -w$ along any path exchanging two particles and for any $w$. 
 We shall first do this in the physically relevant case
 of three dimensions.
For odd dimensions, the configuration space of identical particles is not
orientable. That is to say, the bundle of pseudo-scalars, 
$\bigwedge^{Nd}(T^{*}(\ensuremath{{}^N\ensuremath{\mathbb{R}^{\ensuremath{3}}}}))$, is nontrivial, i.e.\ the selection of a
non-zero global volume form is impossible. The holonomy
representation of this line bundle 
is the alternating character. Complexifying it preserves the 
holonomy and we end up with the Fermi line bundle. 


We now do a different construction which works for all dimensions. 
 Using the  particle number $N$, form the $N^N$ dimensional bundle
$(\ensuremath{\mathbb{C}}^N)^{\otimes \ensuremath{\mathcal{Q}}}$.
Parallel transport along a path is defined in the obvious way  
by using the bijection of the index set induced by the path. 
The 1-dimensional subbundle of totally
antisymmetric elements is the Fermi line for $N$
particles; symbolically, we could write the Fermi line as
$\bigwedge^{\ensuremath{\mathcal{Q}}}(\ensuremath{\mathbb{C}}^N)$. 
To say it another way, the subspace of the fiber that we
want is
the top
dimensional antisymmetric space obtained from  $\ensuremath{\mathbb{C}}^N$ and the
exterior algebra, i.e.\ the determinant line  
$\bigwedge^N (\ensuremath{\mathbb{C}}^N) \subset (\ensuremath{\mathbb{C}}^N)^{\otimes N}$. 
To start to understand this bundle, note that the operation of parallel
transport does map this subbundle to itself. Furthermore, since
parallel transport around a loop permutes the factors, it is easy to
see that the total antisymmetry of the vector 
means the holonomy representation is
equivalent to the alternating character. Explicitly, 
\begin{equation*}
\mathop{\wedge}_{\ensuremath{\boldsymbol{q}} \in \ensuremath{q}} \ensuremath{w}_{\ensuremath{\boldsymbol{q}}} 
\overset{\ensuremath{\Gamma_{\ensuremath{\sigma}}}}{\mapsto} 
\mathop{\wedge}_{\ensuremath{\boldsymbol{q}} \in \ensuremath{q}}{\ensuremath{w}_{\ensuremath{\sigma}\ensuremath{\boldsymbol{q}}}}
 = \det(\ensuremath{\Gamma_{\ensuremath{\sigma}}}) \mathop{\wedge}_{\ensuremath{\boldsymbol{q}} \in \ensuremath{q}} \ensuremath{w}_{\ensuremath{\boldsymbol{q}}}  
= (-1)^{\sigma} \mathop{\wedge}_{\ensuremath{\boldsymbol{q}} \in \ensuremath{q}} \ensuremath{w}_{\ensuremath{\boldsymbol{q}}}
\end{equation*}
where \ensuremath{\Gamma_{\ensuremath{\sigma}}}\ is the holonomy operator associated to the path \ensuremath{\sigma}\
and each $w_\ensuremath{q}$ is an element of $\ensuremath{\mathbb{C}}^N$.


\subsection{Generalization}  

There is an obvious way to generalize the above story. Let $M$ be a
Riemannian manifold and let it be the space in which the particles
move. Furthermore, let $E$ be a cc-Hermitian\ bundle over  $M$; $E$ should be 
the 1-particle value space. Then to generate a theory of identical
particle motion, we do the same as we did before. Let $\ensuremath{\mathcal{Q}} :=
\ensuremath{{}^{N}\! {M}}$.  We define $E^{\otimes \ensuremath{\mathcal{Q}}}$ to have the fiber
$E^{\otimes \ensuremath{q}} \ensuremath{:=} \bigotimes_{\ensuremath{\boldsymbol{q}} \in \ensuremath{q}} E_{\ensuremath{\hat{q}}}$. 
We again define parallel
transport by using the induced paths in $M$ generated by a path in
\ensuremath{\mathcal{Q}}.  The parallel transport of a product element is the  
parallel transport of  the factors in the product  over
the  paths in physical space. In the spin bundle, 
parallel
transport of a factor was trivial.
 We shall refer to
$E^{\otimes \ensuremath{\mathcal{Q}}}$ as the position product bundle of $E$ over $M$
and $\bigoplus_{\ensuremath{\mathcal{Q}}} E$ as the position sum bundle of $E$ over $M$.

Applying the Abelian Quantization Principle\ in any of its forms, 
we again have bosons and fermions as before. We may
also have additional possibilities depending on the dimension of $M$
and its fundamental group. 
 

\section{Beyond the principle?} 

This section describes what happens if we try to go beyond the Abelian Quantization Principle. We
describe this in all three approaches as each has its own flavor to
how the possibility is rejected. Of course, technically we only need
to do this once. Although we are giving this argument in the context
of identical particles with spin, this is really a discussion about
whether the Abelian Quantization Principle\ is formulated in an appropriate level of generality.

\subsection{Beyond \ensuremath{\text{AQP}_{\ensuremath{\widehat{\ensuremath{\mathcal{Q}}}}}}?} \label{s:bprin}


The covering space formulation suggests that the wave function\ on different levels
should differ only by a multiplicative constant. Our motivation was
the realization that this is an invariant condition guaranteeing that
the wave function\ provides a Bohmian\ evolution on the base space. Looking at
Bohm's equation, it is clear that if we replace the constant scalar with a
constant unitary matrix, then we again have a projective Bohmian\
velocity field.  We thus formulate our
periodicity condition\ as  
\begin{equation} \label{e:ansatz}
\ensuremath{\psi}(\ensuremath{\sigma} \ensuremath{\hat{q}}) = \ensuremath{\Gamma_{\ensuremath{\sigma}}}(\ensuremath{\hat{q}}) \ensuremath{\psi}(\ensuremath{\hat{q}})
\end{equation}
where \ensuremath{\Gamma_{\ensuremath{\sigma}}}\ is locally constant.\footnote{
\emph{Local constancy} means that we can obtain \ensuremath{\Gamma_{\ensuremath{\sigma}}}\ from parallel
transporting it along a curve; equivalently, its covariant derivative
satisfies 
\begin{equation*}
\nabla \ensuremath{\Gamma_{\ensuremath{\sigma}}} = 0.
\end{equation*}
Although it makes no sense to ask if it is globally constant since
the vector fibers themselves depend on \ensuremath{\hat{q}}, it does make sense to
ask if $\ensuremath{\Gamma_{\ensuremath{\ensuremath{\sigma}_{1}}}} (\ensuremath{\hat{q}}) = \ensuremath{\Gamma_{\ensuremath{\ensuremath{\sigma}_{1}}}} (\ensuremath{\ensuremath{\sigma}_{2}}
\ensuremath{\hat{q}})$. This equation does not hold; instead we have to apply
permutations:  $\ensuremath{\Gamma_{\ensuremath{\ensuremath{\sigma}_{1}}}} (\ensuremath{\ensuremath{\sigma}_{2}} \ensuremath{\hat{q}}) = R_{\ensuremath{\ensuremath{\sigma}_{2}}} \ensuremath{\Gamma_{\ensuremath{\ensuremath{\sigma}_{1}}}} (
\ensuremath{\hat{q}}) R_{\ensuremath{\ensuremath{\sigma}_{2}}}^{-1}$ where $R_{\ensuremath{\ensuremath{\sigma}_{2}}}$ is the appropriate
permutation matrix representing the parallel transport operator for
spinors. 
}
The 
condition of invariance under
Pauli's equation is a very strong condition; whether \ensuremath{\Gamma_{\ensuremath{\sigma}}}\ is
unitary or not, invariance demands that \ensuremath{\Gamma_{\ensuremath{\sigma}}}\ is a multiple of the
identity. This is what we shall now explain. 

Invariance means that if an initial $\ensuremath{\psi}_0$ satisfies
(\ref{e:ansatz}), then its solution will satisfy (\ref{e:ansatz}) for
all later times. Let $\ensuremath{\psi}_t$  be the solution to Pauli's equation\ with initial
condition $\ensuremath{\psi}_0$ and assume that (\ref{e:ansatz}) holds for all
time. Using Pauli's equation, we  
 compute
\begin{equation*}
  H \ensuremath{\Gamma_{\ensuremath{\sigma}}} \ensuremath{\psi}_t 
= H (\ensuremath{\psi}_t \circ \ensuremath{\sigma}) 
= i \ensuremath{\hbar} \ensuremath{\frac{\partial {(\ensuremath{\psi}_t \circ \ensuremath{\sigma})}}{\partial {t}}}  
= i \ensuremath{\hbar} \ensuremath{\frac{\partial {(\ensuremath{\Gamma_{\ensuremath{\sigma}}} \ensuremath{\psi}_t)}}{\partial {t}}} 
= i \ensuremath{\hbar} \ensuremath{\Gamma_{\ensuremath{\sigma}}} \ensuremath{\frac{\partial {\ensuremath{\psi}_t}}{\partial {t}}} 
= \ensuremath{\Gamma_{\ensuremath{\sigma}}} H \ensuremath{\psi}_t;
\end{equation*}
of critical importance to this computation is the easy fact that if
$\ensuremath{\psi}_t$ is any solution of Pauli's equation, then $\ensuremath{\psi}_t \circ \ensuremath{\sigma}$ is a solution as
well.\footnote{This follows from the fact that composing \ensuremath{\psi}\ with
\ensuremath{\sigma}\ does not affect the Laplacian and that the spin terms for our
bundle are equal at different points along the same covering fiber. In
the standard story, a permutation in spin space would need to be
invoked.} 
We thus have the fundamental relation
\begin{equation} \label{e:commu}
[H, \ensuremath{\Gamma_{\ensuremath{\sigma}}}] \ensuremath{\psi}_t  = 0.  
\end{equation}


By local constancy, \ensuremath{\Gamma_{\ensuremath{\sigma}}}\ commutes with the
Laplacian.
We are left with 
\begin{equation} \label{e:spincom}
[\sum_{\ensuremath{\boldsymbol{q}} \in \ensuremath{\hat{q}}} \ensuremath{(\ensuremath{\boldsymbol{S}}^{\ensuremath{\boldsymbol{q}}} \cdot
{B}_{{\ensuremath{\boldsymbol{q}}}})}, \ensuremath{\Gamma_{\ensuremath{\sigma}}}] \ensuremath{\psi} (\ensuremath{\hat{q}}) = 0.
\end{equation}
To use this equation effectively, we fix a point, $\ensuremath{\hat{q}}$, and
parallel transport our objects from every $\ensuremath{\hat{q}}'$ to \ensuremath{\hat{q}}. Parallel
transport preserves (\ref{e:spincom}). Since \ensuremath{\Gamma_{\ensuremath{\sigma}}}\ is parallel,
we now have a family of equations constraining $\ensuremath{\Gamma_{\ensuremath{\sigma}}}(\ensuremath{\hat{q}})$. For a generic
Hamiltonian, the $\boldsymbol{B}$ field will vary from point to point and we
are led to the conclusion that the commutation relation holds for all
spin terms. Furthermore, if \ensuremath{\psi}\ is sufficiently generic, then we
would expect the commutation relations to hold for all spinors in the
spin space at \ensuremath{\hat{q}}. We can then say that \ensuremath{\Gamma_{\ensuremath{\sigma}}}\ must commute with 
the $3N$ 
element set $\bigcup_{ \ensuremath{\boldsymbol{q}} \in \ensuremath{\hat{q}}
} \{\ensuremath{{S_{x}}}^{\ensuremath{\boldsymbol{q}}}, \ensuremath{{S_{y}}}^{\ensuremath{\boldsymbol{q}}}, \ensuremath{{S_{z}}}^{\ensuremath{\boldsymbol{q}}}\}$;
this set is a set of generators for an irreducible
representation.\footnote{
\ensuremath{{\ensuremath{W}}^{{\otimes}{\ensuremath{q}}}}\ is an irreducible representation space of \ensuremath{\mathop{\times}_{\ensuremath{\boldsymbol{q}} \in \ensuremath{q}} \ensuremath{SU(2)}}, the 
\ensuremath{q}-indexed group product of \ensuremath{SU(2)}. This follows because 
the single-particle
spin space, \ensuremath{W}, is an irreducible representation space of
\ensuremath{SU(2)}; tensor products of representation spaces naturally represent
the group product; and, by Proposition 4.14 on page 82 in \cite{BrDi},
such product representations are irreducible if the factors themselves are
irreducible. The proof in that text uses the theory of
characters of a representation. 
To avoid confusion, please note that we are not discussing the
natural representation of \ensuremath{SU(2)}\ on 
\ensuremath{{\ensuremath{W}}^{{\otimes}{\ensuremath{q}}}}\ which is a reducible representation. 
}
Thus, by
Schur's lemma, \ensuremath{\Gamma_{\ensuremath{\sigma}}}\ must be a multiple of the identity. Again, by
local constancy, 
 the scalar
multiple is a global constant.  We have the Bose-Fermi
alternative and we find that we failed to 
 go beyond the principle. 

\subsection{Beyond \ensuremath{\text{AQP}_{H}}?}

The cut approach goes in much the same way. The difference is that in
the very formulation, we are looking for a domain of
self-adjointness. Just being symmetric implies, 
roughly speaking, that the boundary
conditions need to commute with the Hamiltonian. Thus, we immediately
have the commutation relation between the representation of the
fundamental group and the spin representation leading again to the
fundamental group's representation being essentially one-dimensional.

\subsection{Beyond \ensuremath{\text{AQP}_{\nabla}}?}

As for going beyond \ensuremath{\text{AQP}_{\nabla}}, we need to ask a different
question. Our original question was modest; we only wanted to find a
method for constructing different theories based on a given bundle. We
formulated a principle as to how to do this and we established that it
was equivalent to what we had done in the covering space approach. But
we could also try to find all bundles
that are locally equivalent but
globally inequivalent. As an example, we might know that the theory we
want has the property that the associated bundle $E$ is endowed with a
flat connection. Then the question reduces to characterizing all flat
bundles of a given dimension.
We shall give the well-known answer to this question, and
then explain how one might exclude these possibilities based on other
constraints. In particular, the presence of the spin structures
prevents us from considering all flat bundles
 over \ensuremath{{}^N\ensuremath{\mathbb{R}^{\ensuremath{3}}}}. 


We wish to 
characterize all  
 flat bundles over
\ensuremath{\mathcal{Q}}. Each flat bundle over
\ensuremath{\mathcal{Q}}\ corresponds to a different 
Bohmian theory; locally each flat bundle is equivalent, but globally
these are inequivalent flat  bundles and distinguished by their
holonomy representations.
If the flat bundle has a parallel inner
product,  then the representation is unitary under that inner product.  
After taking into account the appropriate equivalence
relations,\footnote{
Two flat bundles, $E$ and $F$  over \ensuremath{\mathcal{Q}}, are equivalent if there
is a smooth 
diffeomorphism $f$ between 
the two bundles such that $f(\ensuremath{q}, \cdot)$ is an isomorphism from
$E_{\ensuremath{q}}$ onto $F_{\ensuremath{q}}$ for
each $\ensuremath{q} \in \ensuremath{\mathcal{Q}}$ and  the pull-back of the
connection on $F$ agrees
with the connection on $E$. Two representations are equivalent if
there is an isomorphism between the vector spaces which induces an
isomorphism of the representations; that is to say, if $f: V \to W$ is
an isomorphism, $\ensuremath{\sigma} \mapsto A_{\ensuremath{\sigma}}$ is a representation of the
group $G$ on $V$, and $\ensuremath{\sigma} \mapsto B_{\ensuremath{\sigma}}$ is a representation of the
group $G$ on $W$, then $f$ demonstrates the equivalence  if and
only if $f A_\ensuremath{\sigma} f^{-1} = B_{\ensuremath{\sigma}}$. It is an easy fact that if
the representations formed by the parallel transport operators on $E$
and $F$ are
equivalent, then $E$ and $F$ are equivalent as flat bundles. To see
this, simply take the equivalence at one point, parallel
transport it around, and use the representation properties to
demonstrate that this is a consistent procedure, i.e.\ the parallel
transported isomorphism is the same as the original one. The character
of the representation, defined to be the mapping from the group to the
trace of its representative, is in bijective correspondence with the
equivalence classes of the representations. 
}
it is a fact
that for each unitary representation of \ensuremath{\pi_1 ({\ensuremath{\mathcal{Q}}})} acting on $W$, there is
exactly one flat 
bundle with a parallel inner product.\footnote{One way to construct
them is to appeal to the covering space. Let $\ensuremath{\sigma} \mapsto \ensuremath{\Gamma_{\ensuremath{\sigma}}}$ be the
corresponding representation of the covering group. Then define the
equivalence relation $(\ensuremath{\hat{q}}, v) \ensuremath{\sim}
   (\ensuremath{\sigma} \ensuremath{\hat{q}}, \ensuremath{\Gamma_{\ensuremath{\sigma}}} v)$. The corresponding bundle and induced
parallel transport is the correct flat bundle.}
 For each flat bundle,
we have a Bohmian theory; it is a standard fact that equivalent flat bundles
provide equivalent theories.\footnote{Using a connection-equivalence
bundle map, say $\Phi: E \to F$, we may map wave functions to wave functions and note that since
$\nabla^E \Phi \ensuremath{\psi} = \Phi \nabla^F \ensuremath{\psi}$ and $\Phi_{\ensuremath{q}}$ is
unitary, we have the same Bohmian velocities. As $\Phi$ also suitably commutes
with the Hamiltonians, we have the same Bohmian evolution. 
} Thus, for flat bundles whose potentials are scalar
multiples of the identity, we have a different theory for each unitary
representation of \ensuremath{\pi_1 ({\ensuremath{\mathcal{Q}}})}.

In the case of identical particles with spin,
 our potentials are not just scalar multiples of the identity. We
want a bundle which can support the spin terms. 
Our demand is that
that the spin operators should locally form a
parallel set of operators and be appropriately linked to the physical
positions in space.  Specifically, we demand that for each element
$\ensuremath{\boldsymbol{q}} \in  \ensuremath{q} 
\in \ensuremath{\mathcal{Q}}$, there
is an associated  spin vector, $\ensuremath{\boldsymbol{S}}_{\ensuremath{\boldsymbol{q}}}$.  As we parallel
transport these 
operators along a 
path $\alpha$ in \ensuremath{\mathcal{Q}}, the spin vector associated to $\ensuremath{\boldsymbol{q}}$ should be
mapped into the spin vector associated with
$\alpha_{\ensuremath{\boldsymbol{q}}}(t)$, the point in physical space at time $t$ that 
$\ensuremath{\boldsymbol{q}}$ is mapped to via $\alpha$.\footnote{Recall that a path $\alpha
: [0,1] \to  
\ensuremath{\mathcal{Q}}$ is a set of $N$ trajectories in physical space. We denote the
path starting at \ensuremath{\boldsymbol{q}}\ by $\alpha_{\ensuremath{\boldsymbol{q}}} : [0,1] \to \ensuremath{\mathbb{R}^{\ensuremath{3}}}$.}  
Around a representative of a nontrivial element in \ensuremath{\pi_1 ({\ensuremath{\mathcal{Q}}, \ensuremath{q}})},
the spin vectors are permuted.\footnote{ 
This implies that the spin operators are not sections of the
endomorphism bundle. Only the set of spin vectors is globally defined.
We could define a local parallel section for each spin operator, but
these cannot be extended globally on \ensuremath{\mathcal{Q}}. 
}
The holonomy operators are responsible for implementing this permutation. 
By using Schur's lemma appropriately, we can argue that the holonomy
representation  must either be the permutation representation or the
alternating permutation 
representation.\footnote{If we know the action of a matrix on all of
 the elements of an irreducible representation, then we know the
 matrix up to a scalar multiple. Since the action is that of a
 permutation, then the holonomy must be of the form of a character
 tensored with the permutation representation.}


 Thus, the only flat bundles capable of supporting the
spin terms are the Bose bundles and the Fermi bundles.
 The argument given above is the argument given in \cite{LM77}.
Thus, if we demand flatness and the presence of spin, we conclude that
our principle is the most generic that we can expect. 

What if we remove the constraint of flatness? 
Trying to characterize all bundles of a certain dimension with a
prescribed curvature, is not, to our knowledge, 
an easy task. Our general principle allows us to construct new bundles
from a given bundle, but to find all such bundles is an entirely
different question. With the spin terms, we can at least substantially
reduce the question. Indeed, we still demand the spin operators to be
parallel and to be permuted upon transporting around any loop. 
Thus, along any path, Schur's lemma tells us that 
parallel transport along a closed
curve is implemented by scalar multiples of permutation matrices. The
permutation depends only on the homotopy class of the curve while the
scalar will generally vary over the homotopy class. It then seems
plausible to say that all bundles with the spin terms may be formed by
tensoring the spin bundle with a line bundle. The question of
classification then reduces to classifying all line bundles over
\ensuremath{\mathcal{Q}}.\footnote{From a discussion with R. Tumulka, it seemed that
the key question is whether there are closed, but not exact
2-forms. This is because the curvature tensor for a line bundle is
a closed 2-form. Thus, if all closed 2-forms are exact, then given a
prescribed curvature, there is a 1-form which may be subtracted from
the connection to achieve a flat connection; since flat connections
have been characterized by $\pi_1$, this would conclude the argument. 
But closed
2-forms being exact is implied by the triviality of the second
homotopy class, $\pi_2(M)$. 
In other words, if all continuous maps of $S^2$ into
$M$ are 
contractible in the space, then 2-forms are exact. 
Except for $d=3$,  $\pi_2(\ensuremath{\ensuremath{\mathbb{R}^{d,N}}_{\neq}})$ seems to be
trivial. Indeed, pick two particles and place a sphere between two of
them. This sphere can be contracted iff $d \geq 4$; this follows by
viewing one particle as the origin and demanding that the contracting
 sphere never intersects the origin. We would suspect that, for $d
\geq 4$, $\pi_2(\ensuremath{{}^N\ensuremath{\mathbb{R}^{\ensuremath{3}}}})$ is finite and therefore its cohomology is trivial. 
Whether or not the exceptional case $d=3$ is an exception is
not known to us. We conjecture that, just as anyons arise in two
dimensions, in three dimensions we do have such 2-forms and quite
possibly such bundles. 
}


\section{Below the principle} \label{iss:beypri}


Just as it is possible in special circumstances to formulate theories
that go beyond the Abelian Quantization Principle, it is also possible to be in a situation in
which the different characters give rise to equivalent theories. We
have put this material here since the example is best understood at
this point. 


This bundle was suggested
by Federico Bonetto in the context of a theory of mixed bosons and
fermions; we discuss this here only in the scalar case. One approach to a
system in which we have $m$ bosons and  $n$ fermions is to 
to choose the configuration space to be the 
Cartesian product ${}^{m} \ensuremath{\mathbb{R}^{3}} \times {}^{n} \ensuremath{\mathbb{R}^{3}} $. The
 bundle over that space would then be the ordinary tensor product of the 
 the respective Bose and Fermi lines. Another possibility is to 
take the special $3k$-forms on \ensuremath{{}^N\ensuremath{\mathbb{R}^{3}}}\ which
respect the particle identity; in local coordinates, we have 
 $d\ensuremath{\boldsymbol{q}}_{i_1} \wedge  \cdots \wedge d\ensuremath{\boldsymbol{q}}_{i_k}$.


Such a $3k$-form could be thought to represent
$k$-fermions and $N-k$ bosons. What happens if we tensor this bundle
with the Fermi line? As one would naively expect, this switches the
fermions and bosons, i.e., it maps the $3k$-forms into
$3(N-k)$-forms. Thus, 
the bundle of all special forms is equivalent to itself tensored with
the Fermi line. In order for the two Bohmian theories to be
equivalent, we would actually need the Hamiltonian to be the same
under the equivalence;  this seems to place some restrictions on the
Hamiltonian. 
The interested reader can find more details about the interplay of
different types of particles and \ensuremath{{}^N\ensuremath{\mathbb{R}^{3}}}\ in
\cite{GTu03} and Section {\ref{s:dpart}}. 


For a simpler counterexample,  take 
the bundle $E$ to be the direct sum  (Whitney sum) of
the Fermi line, $F$,  and the Bose line, $B$,  over \ensuremath{{}^N\ensuremath{\mathbb{R}^{3}}}. Then the principle
 suggests that to find other 
theories, we should tensor the Fermi line, $F$,
with $E$ to obtain a new bundle $G$. 
But in fact, we have 
\[
G \ensuremath{:=} F \otimes (F \oplus B) \ensuremath{\cong} F \otimes F \oplus F \otimes B
\ensuremath{\cong} B \oplus F \ensuremath{\cong} E . 
\]
Assume that we had a potential $V= \bigl( \begin{smallmatrix} a &b \\
c&d \end{smallmatrix} \bigr)$ acting on $E$. If one computes, one
finds that $\ensuremath{\text{Id}} \otimes V \mapsto \bigl( \begin{smallmatrix} d &c \\
b&a \end{smallmatrix} \bigr)$ under these isomorphisms. 
As derivatives are unaffected by these
isomorphisms,  we shall have equivalent theories iff $a=d$ and $b
=c$. 


Note
that what happens in these case is  the vanishing of 
the trace of the holonomy operators of $E$
representing  the paths whose holonomy operators in $F$  are
 non-trivial.
 Since the character of a
representation determines a representation up to
isomorphism 
and the
character of a tensor product is the product of the
characters,\footnote{See \cite{BrDi}. Recall that the character of a 
representation is the mapping from the group into the complex numbers
defined by taking the trace of the representatives.} we have that
$G$'s holonomy representation is
 equivalent to $E$'s holonomy representation.
It seems that if we avoid this kind of degenerate situation, 
then we will obtain inequivalent
theories. 


\section{Bundle triviality}

A natural question to ask about the Bose bundles and Fermi bundles is
whether these bundles are trivial or not. 
They are certainly not trivial flat
bundles. But a bundle can have a non-trivial flat connection and still
be trivial as a bundle, 
as is the case for the complex line over the circle. It
turns out that some of the bundles are trivial and some are not. Given
that information, it might then be reasonable to suggest that the
spin-statistics relation would be correlated to the triviality of the
bundle. More precisely, the conjecture would be that 
 the Bose bundles for integral spin
 would be trivial while the Fermi bundles would
be trivial for half-integral spin. This conjecture is false. In
fact, there is no correlation. Here is the table of relations for
$N=2$, $d =3$, spin $s$:
\begin{center}
\begin{tabular}{|l|c|c|c|c|}
\hline
 & $s=2l$  &  $s=2l+\frac{1}{2}$ &  $s=2l+1$ &  $s=2l+\frac{3}{2}$\\
\hline
Bose & trivial & nontrivial & nontrivial & trivial \\
Fermi  & nontrivial & nontrivial & trivial & trivial\\
\hline
\end{tabular}
\end{center}
for $l \in \mathbb{N}$. 
For $N > 2$, the Bose bundle  is nontrivial for odd integral spin while 
the Fermi bundle is nontrivial for even integral spin. We do not know the 
triviality or nontriviality of the other bundles except for the Bose
line which is clearly trivial. We give here a
 brief outline of the reason for these
results. In Appendix \ref{s:triv}, we give a more complete account of these
results.


We start with some facts about line bundles over \ensuremath{{}^N\ensuremath{\mathbb{R}^{\ensuremath{d}}}}. First, a 
flat line bundle with trivial holonomy is a trivial bundle. To
trivialize, simply parallel transport a (1-element) basis to every
point of the manifold. Since the
holonomy is trivial, this is a well-defined operation. Second, if a
flat line bundle has non-trivial holonomy, over \ensuremath{{}^N\ensuremath{\mathbb{R}^{\ensuremath{d}}}}, then it is
non-trivial. As there are only two characters of the fundamental
group, this is a statement directly about the Fermi line. The
proof of this assertion lies in the following idea.
The difference between
two connections is a 1-form and for both of them to be flat, the
1-form must be closed. Furthermore, to have a nontrivial holonomy, the
1-form must be non-exact. Since all closed 1-forms on \ensuremath{\mathcal{Q}}\ are exact, we
are done. Thus the Fermi line is nontrivial. Furthermore, the
configuration space's top form bundle has the holonomy of the
alternating character iff $d$ is odd. This follows by considering the
wedge product of the basis elements and noting that the parallel transport
along a loop transposing two particles involves the exchange of 
 $d$-pairs of vectors. Thus, \ensuremath{{}^N\ensuremath{\mathbb{R}^{\ensuremath{d}}}}\ is orientable iff $d$ is even.

A useful tool in establishing the nontriviality of a vector bundle is
its determinant line bundle. If $E$ is a vector bundle of dimension $l$,
then its determinant line is $\det(E) \ensuremath{:=} \bigwedge^l (E) = (E^{\otimes
l})_{Alt}$, the totally antisymmetric, 1-dimensional subbundle of the
$l$-fold tensor product of $E$ with itself. If $E$ is a trivial
bundle, then its determinant line is trivial. The converse fails to
hold as the case of the nontrivial tangent bundle to the (orientable)
2-sphere demonstrates. 
We also have that the parallel transport operator on $E$ induces a
parallel transport 
operator on $\det(E)$. In fact, for a given path, the 
holonomy operator for that path in $\det(E)$ is multiplication by the
determinant of the corresponding holonomy operator acting on $E$. This
is related to why it is called the determinant line bundle. Furthermore, if $E$
is flat, then $\det(E)$ is flat. Hence our plan is to compute the
determinant line bundle of our bundles. This is equivalent to computing
the determinant of the matrices in the permutation and alternating
 permutation representation. 

Given any of the bundles, our candidate permutation matrix will be the
one corresponding to an exchange of two particles. As is easy to
see,\footnote{For $N=2$, decompose the tensor product into symmetric
and antisymmetric elements. For $N>2$, consider a matrix representing
a transposition. Then again decompose it into the symmetric and
antisymmetric parts with respect to those two factors.} 
we find that the determinant of the matrix for the Bose bundle is 
$-1$ raised to the power $\binom{2s+1}{2}(2s+1)^{N-2}$ while the
determinant for the 
matrix for the Fermi bundle is $-1$ raised to the power $(2s+1 +
\binom{2s+1}{2})(2s+1)^{N-2}$. By computing when the exponents are
odd, we arrive at our claims about the non-triviality. 

Our final effort will be to demonstrate the triviality assertions. To
do this, we need to argue that
the direct sum 
of the Fermi line with itself, over ${}^2\ensuremath{\mathbb{R}^{3}}$, is trivial, i.e.\
$F \oplus F \ensuremath{\simeq} {}^2\ensuremath{\mathbb{R}^{3}}\times \ensuremath{\mathbb{C}}^2 $. 
The following proof was implied to exist in appendix D of \cite{MR97}.  
The idea is to produce a trivialization. We start with writing the
covering space as
\[
 \ensuremath{\widehat{\ensuremath{\mathcal{Q}}}} \ensuremath{:=} \ensuremath{\mathbb{R}^{2,3}}_{\neq} = \ensuremath{\mathbb{R}^{3}} \times \ensuremath{\mathbb{R}^{+}} \times
 S^2,
\]  
i.e.\   center of mass, distance between the particles, and the
direction from particle 1 to particle 2. To create the base space, we
identify antipodal points on $S^2$.
To find a basis for our bundle, it is sufficient to find a global basis of
$S^2 \times \ensuremath{\mathbb{C}}^2$ such that both sections of the basis satisfy the
alternating periodicity condition. 
Once we have such sections, we then extend them to \ensuremath{\widehat{\ensuremath{\mathcal{Q}}}}\ by
defining them to be constant in the other coordinates.
To obtain these sections over $S^2$,
 we parameterize the sphere with spherical coordinates $(\theta, \phi)$:
 $0
\leq \theta \leq 2\pi$ and $0 
\leq \phi \leq \pi$. Then the following 
smooth, linearly independent sections 
\[
\ensuremath{\psi}_1 (\theta, \phi)= 
\begin{pmatrix}
\cos\phi\\
e^{i\theta}\sin\phi
\end{pmatrix}
\quad
\ensuremath{\psi}_2 (\theta, \phi) = 
\begin{pmatrix} 
e^{-i\theta} \sin\phi\\
-\cos\phi
\end{pmatrix}
\]
satisfy   
\[
\ensuremath{\psi}_i(\theta + \pi, \pi - \phi) = - \ensuremath{\psi}_i(\theta, \phi).
\] 
This is all we needed. This  comes from considerations in the
spin-$\frac{1}{2}$ representation of $SO(3)$. Namely, the matrices 
$A(\theta,
\phi)$ with columns $\ensuremath{\psi}_1$ and $\ensuremath{\psi}_2$, represent a
rotation of angle $\pi$ about the axis specified by $(\theta, \phi)$.
The essential fact is that we have a double-valued representation,
i.e.\  $A(\theta + \pi, \pi - \phi) = -A(\theta, \phi)$.

The following works for $N=2$. 
The idea is to decompose each bundle into a sum of Bose
and Fermi lines. 
We then take pairs of  Fermi lines and trivialize them. 
If there are an even number of Fermi lines, then we have
trivialized the whole bundle as the sum of trivial bundles is
trivial. 
 An odd number implies that it is
nontrivial, a fact we have already established. 
We start with the Bose bundles.
Since we have only two particles,
the bundle is the sum of a symmetric part and an antisymmetric
part. The symmetric subbundle is trivial (a sum of Bose lines), 
while the antisymmetric part
is a sum of Fermi lines. Thus, the Bose bundle is trivial if and only
if the number of Fermi lines is even i.e., iff  $\binom{2s+1}{2}$, the
dimension of the  
antisymmetric part, is even. 
As for the
Fermi bundle, its triviality depends upon the parity of $2s+1 +
\binom{2s+1}{2}$, the dimension of the
symmetric part, since upon tensoring with a Fermi line, the Fermi lines
become Bose lines and the Bose lines become Fermi lines. 
The results follow.


\section{Another viewpoint for identical particles} \label{s:view}


We have described 
how a theory about identical particles should be formed in Bohmian mechanics. In
this section, we shall pursue a different, very speculative, idea. The
idea suggests
how Bohmian mechanics\ can choose bosons and fermions even if we start on the
wrong space. Furthermore, it demonstrates 
the difficulties 
in 
 satisfying the intuition that permutation symmetry leads to
only two classes of wave functions. It is also useful in revealing a certain
strength of Bohmian mechanics. But we are not suggesting that this is the
appropriate story for identical particles; that story has already been
presented. 
Unless otherwise stated, we will be discussing complex-valued
wave functions even though spin does not substantially change the discussion. 
The wave functions will be functions on \ensuremath{\mathbb{R}^{\ensuremath{d} N}}\ and the configurations will
be elements of \ensuremath{\mathbb{R}^{\ensuremath{d} N}}. 

We have a system of labelled particles. 
In this context,  a system of identical particles 
requires two things. The first is to define an action of the
permutations on the state of the system. Once one has that, then the
particles are identical iff the action of the permutation group
commutes with the evolution of the system. In other words, if $A_t$ is
the evolution operator for the 
system, $\chi$
is any initial condition, and $\sigma$ is a permutation,
then we require that $\sigma(A_t(\chi)) = A_t(\sigma\chi)$. 

The second condition is straightforward to understand, but the first
condition can be tricky.  The action of a permutation on the particle
positions is, of course, to exchange the positions. One must decide
how the other objects
in the theory should change. For this, we have two classes of
objects. Roughly, objects considered as laws follow the particles while 
state objects follow the positions.
As that was a bit vague, let us look at classical mechanics.
In classical
mechanics, the state of the system is given by the position and
velocity of the particles. When a permutation acts on the sate, the
tangent vector associated with the given position stays associated
with the same position. But both the position and tangent vector are
now associated to a different particle. In contrast, 
the law is given by the force which is a function of the state of the
system. 
We do not redefine the force in order
to undo the permutation; rather, we take the state as it is after the
action of the permutation and evaluate the force at that position. 
We would therefore generally expect the
initial force vectors to have changed when we act upon the initial
state by a permutation. We can then see that the second condition will be
satisfied in classical mechanics iff the force is invariant under
permutations, i.e.\ we need, say for two particles, that 
$F(\ensuremath{\boldsymbol{q}}_a, \dot{\ensuremath{\boldsymbol{q}}_a}, \ensuremath{\boldsymbol{q}}_b, \dot{\ensuremath{\boldsymbol{q}}_b})
= F(\ensuremath{\boldsymbol{q}}_b, \dot{\ensuremath{\boldsymbol{q}}_b}, \ensuremath{\boldsymbol{q}}_a, \dot{\ensuremath{\boldsymbol{q}}_a}).$

How do permutations act in Bohmian mechanics? That is the question which is central
to this section. There are three objects in the theory: $Q = 
(\boldsymbol Q_1, \ldots, \boldsymbol Q_N) \in \ensuremath{\mathbb{R}^{\ensuremath{d} N}}$, \ensuremath{\psi},
$H$. The status of $Q$ is clear and permutations act on it by
permuting the positions, i.e.\ $( \boldsymbol Q_1, \ldots, \boldsymbol
Q_N) \mapsto \boldsymbol Q_{\sigma 1}, \ldots, \boldsymbol Q_{\sigma
N}$. The Hamiltonian, being a law, is unchanged as an object acting on
the wave functions. The question is how \ensuremath{\psi}\ should change. 
The standard view of Bohmian mechanics, as stated in Section {\ref{s:sim}}, treats
\ensuremath{\psi}\ as part of the state of the system. Thus, we have the freedom,
and the obligation, to have the permutations act on it. Since  \ensuremath{\psi}\
generates the velocity field, we want the permutation action to be
such that the velocity vector at $\ensuremath{\boldsymbol{q}}$ is unchanged when we exchange
the particles. That is to say, we want 
$ v^{\sigma \ensuremath{\psi}}_{\sigma^{-1} i}  (\ensuremath{\boldsymbol{q}}_{\sigma 1}, \ldots
\ensuremath{\boldsymbol{q}}_{\sigma N})  
= v^{\ensuremath{\psi}}_i (\ensuremath{\boldsymbol{q}}_1, \ldots \ensuremath{\boldsymbol{q}}_N). 
$ To understand the formula, 
note that if $k = \sigma^{-1} i$, then 
$\ensuremath{\boldsymbol{q}}_i = \ensuremath{\boldsymbol{q}}_{\sigma k} $ implying that the $\ensuremath{{k}^{th}}$
particle has position $\ensuremath{\boldsymbol{q}}_i$ after the action of the
permutation; hence, the velocity vector of the $\ensuremath{{k}^{th}}$ particle, after the
permutation, should be the same as the velocity vector of the
$\ensuremath{{i}^{th}}$ particle before the permutation. 
As can be readily checked, if $\sigma \ensuremath{\psi} \ensuremath{:=} \ensuremath{\psi} \circ
\sigma^{-1}$, then the above equality holds. Defining the permutation
action on wave functions in this manner, we see that if $H$ is invariant under
permutations, then the system satisfies the conditions for it to be a
system of identical particles. We do not arrive at any constraint on
the wave function. From this viewpoint, the experimental evidence of particles
being either bosons or fermions must be incorporated as an additional
postulate and is not a result of the particles being identical. 
Succinctly, the problem is that the
wave function\ changes covariantly under symmetries.\footnote{As far as we can
tell, this is also a problem in standard quantum mechanics as
evidenced by the symmetrization postulate. In the standard story,
the wave function\ should change in the same way, resulting in the
probability density changing in just the right way so as to be
insensitive to a relabelling.} 

The remedy of this problem, motivated by other reasons, is to consider
the wave function\ as part of the law. 
It is our contention that the
wave function\ should be thought of as an analog of the Hamiltonian; that is to say,
the wave function\ defines the law of motion and can be given before the
evolution of the particles is considered.\footnote{
This is true for isolated Bohmian\ systems or, more appropriately, for
the universal wave function. For systems that are
interacting  
with their environment, the system's wave function, defined to be the composite
wave function\ (system plus environment)
evaluated at the actual configuration of the environment, can evolve
in a way different from the Schr\"odinger\ evolution. This is the case  when
a wave function\ collapses \cite{DGZ92}. Effectively, a subsystem is governed by two
complementary evolutions: the Schr\"odinger\ evolution
and an evolution based on
the environment's configuration changing. 
}
For details on the wave function\ as a law, please see \cite{DGZ97}. This
viewpoint implies that the velocity 
field  is not a consequence of the dynamics, but rather defines the
dynamics in the same way as the force does in classical mechanics. 
The only initial conditions are the initial positions of the particles. 
Accepting our assertion leads to the conclusion that the velocity
field should be appropriately invariant under permutations of the
configurations. This eventually leads to concluding that we have either
bosons or fermions.


A natural challenge to our invoking ``the wave function\ is a law'' argument is
why we would still adhere to the usual policy of changing  the wave function\ under a
rotation or gauge transformation, but not under a permutation
symmetry. 
The answer lies with the conditional wave function. If we have a
wave function\ governing a system and its environment, we can write it as a
function of two variables; namely, the wave function\ is a function of the
configuration of the system and the configuration of the environment:
$\ensuremath{\psi}(\text{sys, env})$. We 
then define the conditional wave function\ of the system to be $\ensuremath{\psi}(\cdot,
\text{ENV})$ where ENV stands for the actual configuration of the
environment. In certain situations, the conditional wave function\ will satisfy
its own Schr\"odinger's equation\ and give rise to a Bohmian\ system whose trajectories for
the system will agree with the actual Bohmian\ trajectories. See
\cite{DGZ92} for more details about the conditional wave function. 


If we have a fundamental
symmetry of the universal wave function, then there is, in general, no reason to
assume that symmetry for a conditional wave function. For example, consider
rotations. If our universal wave function\ is invariant under rotations of
physical space, then rotating our initial universal configuration
produces an equivalent motion. But now let us consider a
smaller system. If we rotate the system's configuration, but not the
environment's configuration, then there is no reason to assume
invariance. Invariance only holds when both system and environment are
transformed, not when just one of them is transformed. 

Symmetry under exchange of particle positions is quite different. This
symmetry is inherited by conditional wave functions. This is obvious; the
exchange symmetries work for any subcollection of identical
particles. Thus, if we believe that the particles are fundamentally
identical, then we may assume this property for all systems. It is
relevant to our earthly domain of physics and it is why our argument
in this section ultimately works. 

It should be quite clear that this is a Bohmian argument. In sharp
contrast to 
 the other approaches presented in this chapter, we can see no way of
doing this argument in standard quantum mechanics.


\section{Classical Mechanics and identical particles} \label{s:cm} 


It is interesting to contrast the Bohmian\ story of identical
particles with that of the classical story. For identical, classical particles,
we would again use the natural configuration space. The initial
conditions are the positions and velocities of the particles. The law
is specified by an acceleration field on configuration space. 

If we demand that  the acceleration field  comes from a potential,
we might expect the covering space to be relevant. But lifting to
\ensuremath{\ensuremath{\mathbb{R}^{d,N}}_{\neq}}\ does not achieve any greater generality. To explain why  \ensuremath{\ensuremath{\mathbb{R}^{d,N}}_{\neq}}\
is not useful, 
consider the periodicity condition\ for the proposed potential $V$:  
\begin{equation*}
V(\ensuremath{\sigma} \ensuremath{\hat{q}}) = V(\ensuremath{\hat{q}}) + \ensuremath{\gamma_{\ensuremath{\sigma}}}.
\end{equation*}
This is the general condition that will lead to the same gradient at
each level of the covering fiber. 
The map $\ensuremath{\sigma} \mapsto \ensuremath{\gamma_{\ensuremath{\sigma}}}$ forms an additive representation  of the
fundamental group. It is impossible to form an additive
representation, over the reals, 
of a finite group. For example,  consider a
transposition. Writing down the periodicity condition s leads to $2 \ensuremath{\gamma_{\ensuremath{\sigma}}} =
0$. There is only one solution to that equation. 

To put it succinctly, all closed 1-forms on \ensuremath{{}^N\ensuremath{\mathbb{R}^{\ensuremath{d}}}}, $d \geq 3$, are
exact. In classical mechanics, the covering space of identical
particles seems irrelevant except as a convenient setting for analysis.
Incidentally, the natural configuration space is useful in classical
mechanics and statistical mechanics. 
Indeed, it answers the Gibbs paradox of thermodynamics; see
\cite{LM77} for details.

\section{Distinguished particles} \label{s:dpart}

The question arises as to what the proper description of distinguished
particles should be. The typical viewpoint is that the parameters are
attached to the particles. In other words, there are electrons,
protons, etc. with various charges and masses 
and these are different types of particles. 
In a typical approach, one would take the Cartesian product of 
 the natural configuration spaces for the various particles, 
possibly removing coincident
points, and tensor the appropriate bundles together to form the bundle
whose sections are the wave functions. 

In classical mechanics, the Cartesian product seems to be the only
option. We see no other way of giving the points different fundamental
properties. But in  
quantum mechanics, we do have another
option, as hinted at by Bell in \cite{Bell87}. The idea is that the particle
differences are all incorporated into the wave function\ itself. We shall
describe the setup here. See \cite{GTuTaZ03} for more discussion. 

We start with the scalar wave functions and the usual theory on \ensuremath{\ensuremath{\mathbb{R}^{d,N}}_{\neq}}. In
general, the velocity field generated by a wave function\ on \ensuremath{\ensuremath{\mathbb{R}^{d,N}}_{\neq}}\ will have no
particular symmetry properties. The simple idea is to symmetrize the
probability current as well as the probabilties and then define the
Bohmian velocity from that. We then have a well-defined velocity field on 
\ensuremath{{}^N\ensuremath{\mathbb{R}^{\ensuremath{d}}}}. Furthermore, if the wave function\ only has support in one ``level'' of
\ensuremath{\ensuremath{\mathbb{R}^{d,N}}_{\neq}}, then this velocity field gives the usual motion for
distinguished particles.  
In particular, the particles will retain their distinctive
motions. We also note that the corresponding equivariant density for
this new motion is the symmetrized density; one could argue that this is
the only position density that one has access to. Indeed, based on the
position of particles at a single time, it would be hard, if not
impossible, to distinguish an electron from a proton if they were
indeed point particles. It is only
through their paths that we see their differences. 

We can write the above theory in a better way. In fact, on \ensuremath{{}^N\ensuremath{\mathbb{R}^{\ensuremath{d}}}}, we
form the $N!$-dimensional bundle obtained by summing over the 
covering fibers. More precisely, at $\ensuremath{q} \in \ensuremath{{}^N\ensuremath{\mathbb{R}^{\ensuremath{d}}}}$, we form the
bundle $\bigoplus_{\ensuremath{\ensuremath{\rho}}^{-1}(\ensuremath{q})} \ensuremath{\mathbb{C}}$. We have in mind that a
wave function\ $\hat{\ensuremath{\psi}}$ on \ensuremath{\ensuremath{\mathbb{R}^{d,N}}_{\neq}}\ defines $\ensuremath{\psi}(\ensuremath{q}) \ensuremath{:=}
 \frac{1}{N!} \oplus_{\ensuremath{\hat{q}} \in \ensuremath{\ensuremath{\rho}}^{-1}(\ensuremath{q})} \ensuremath{\psi}(\ensuremath{\hat{q}})$.  
Each component satisfies a Schr\"odinger\ equation and the sum over the
probability currents leads to the correct symmetrized Bohmian velocity of
$\hat{\ensuremath{\psi}}$. 

We now give a much more general story about distinguishable
particles. 
Fix $K$ to be the number of different types of particles; we shall use
$k$ to index the different particles.  Let
$(E_k, \nabla_k, m_k, \ensuremath{\boldsymbol{S}}_k)$ be the data for particle type
$k$ and we take $e_k$ to be the dimension of the fiber. 
Then the one-particle bundle is $E \ensuremath{:=} \bigoplus_{k=1}^K E_k$ with
the connection induced from a direct sum, e.g.\ $\nabla (\ensuremath{\psi}_1 \oplus
\ensuremath{\psi}_2) \ensuremath{:=} \nabla_1 \ensuremath{\psi}_1 \oplus \nabla_2 \ensuremath{\psi}_2$. The Schr\"odinger\
Hamiltonian is defined as $(G\otimes M + \sum_k \ensuremath{\boldsymbol{S}}_k ) \circ D^2 \ensuremath{\psi}$
where $M$ is a diagonal matrix with the masses along the
diagonal. In each component of the sum, it will look like the usual Pauli
equation. Bohm's equation is dealt with in the same way, i.e. use $G
\otimes M$ to make the conversion from a 1-form to a tangent vector. 
It is necessary to do this before taking the
inner product. 


For many particles, the configuration space will be \ensuremath{{}^N\ensuremath{\mathbb{R}^{\ensuremath{d}}}}. Let \ensuremath{\mathcal{K}}\
denote the set $\{1, \ldots K\}$. 
Let $\ensuremath{A_{q/K}}$ be the set of partitions of $\ensuremath{q}$ into $K$ 
sets. We shall use \ensuremath{\nu}\ to denote a partition which may also be
viewed as a map $\ensuremath{\nu} : \ensuremath{q} \to \ensuremath{\mathcal{K}}$.  
We denote the set of positions of type $k$, which is the set $\nu^{-1}
(k)$,
 by $\ensuremath{\ensuremath{q}_k^{\nu}}$. 
The bundle
for a theory just about bosons is $E^{\otimes \ensuremath{\mathcal{Q}}}$. The fiber at
$\ensuremath{q}$ may be written as 
\[
E^{\otimes \ensuremath{\mathcal{Q}}}_{\ensuremath{q}} \ensuremath{:=} \otimes_{\ensuremath{\boldsymbol{q}} \in \ensuremath{q}}
E_{\ensuremath{\boldsymbol{q}}} \ensuremath{\cong} \bigoplus_{\ensuremath{\nu} \in \ensuremath{A_{q/K}}} (\bigotimes_{k \in \ensuremath{\mathcal{K}}}
\{E_k^{\otimes 
\ensuremath{\ensuremath{q}_k^{\nu}}}\}).\]   
The
connection is defined in
the same way as the spin bundle's connection: it
is induced by the set bijections which 
paths in \ensuremath{{}^N\ensuremath{\mathbb{R}^{\ensuremath{d}}}}\ provide. Parallel transport in terms of the 
 sum over partitions requires taking into consideration that 
 the bijection also maps one partition to another. 

For a theory involving fermions, we must construct
the appropriate bundles first. 
Given a  subset $r$ of the point $\ensuremath{q} \in \ensuremath{\mathcal{Q}}$, we
define the vector space 
\[
F_r \ensuremath{:=} \mathop{\wedge}_{r} \ensuremath{\mathbb{C}}^{\ensuremath{\lvert {r} \rvert}} \ensuremath{:=}
(\ensuremath{\mathbb{C}}^{\ensuremath{\lvert {r} \rvert}})^{\otimes r}_{Alt}.\] We can therefore define
$\bigwedge^i_{\ensuremath{\mathcal{Q}}}$ 
as the bundle whose fiber is the direct sum of the $F_r$'s
where $\ensuremath{\lvert {r} \rvert} =i$. If we have any path $\sigma$ in \ensuremath{{}^N\ensuremath{\mathbb{R}^{\ensuremath{d}}}}\ which gives a
bijection between $r$ and itself, then we have that parallel transport
around that path will be $(-1)^{{\alpha_r}}
(1)^{{\alpha_{\ensuremath{q} \setminus r}}}$. That is to say, it is the
sign of the permutation induced on the set $r$ by the bijection
$\sigma$. In other words, we have, in a loose sense,  $i$
fermions and $N-i$ bosons. As one can easily see, the subspace 
of the points in $r$ being fermions. 
is not invariant whereas the number of fermions is invariant under
the parallel transport. 
In three dimensions, this is equivalent to the $3i$-form bundle of
Section {\ref{iss:beypri}}.  The
direct sum over $i$ leads to the bundle of all mixed bosons and
fermions; we denote it by $\bigwedge_{\ensuremath{\mathcal{Q}}}$. 

The situation where the first $j$ of the $K$ type of
particles are fermionic and the other types are bosonic may be easily
described using these restricted Fermi lines. 


The appropriate bundle is $J$ and it is defined at a point $\ensuremath{q}$ to
be 
\[
J_{\ensuremath{q}} \ensuremath{:=} \bigoplus_{\ensuremath{\nu} \in \ensuremath{A_{q/K}}} (\{\bigotimes_{k =1}^j 
F_{\ensuremath{\ensuremath{q}_k^{\nu}}} \otimes (E_k)^{\otimes \ensuremath{\ensuremath{q}_k^{\nu}} }\} 
\otimes \{\bigotimes_{k =j+1}^K 
 E_k^{\otimes \ensuremath{\ensuremath{q}_k^{\nu}}}\}).\]
 Its dimension is 
\[
\sum_{\ensuremath{\nu} \in \ensuremath{A_{q/K}}}
\binom{N}{l_1\cdots l_K} \{ \prod_{k \in \ensuremath{\mathcal{K}}} e_k^{l_k}\} = 
(\sum_{k \in \ensuremath{\mathcal{K}} } e_k)^N.
\] 
The evolution 
equations are the same as they always have been. Explicitly, we have
\begin{subequations}
\begin{align} 
\ensuremath{\frac{d{\boldsymbol{Q}_t}}{d{t}}} =&  \hbar  \textrm{Im}
\frac {\ensuremath{({\ensuremath{\psi}_t}, { (G_{\boldsymbol{Q}_t} \otimes M)^{\oplus\otimes} \circ
D_{\boldsymbol{Q}_t} \ensuremath{\psi}_t})}}
{\ensuremath{({\ensuremath{\psi}_t}, {\ensuremath{\psi}_t})}} (Q_t)
  \quad \text{for each $\boldsymbol{Q}_t \in Q_t$,} \\ 
i \hbar \ensuremath{\frac{\partial {\ensuremath{\psi}_t}}{\partial {t}}}  (\ensuremath{q}) =&  - \frac{\hbar^2}{2} \sum\limits_{ \ensuremath{\boldsymbol{q}} \in \ensuremath{q}}
 \{ (G_{\ensuremath{\boldsymbol{q}}} \otimes M + \sum_{k \in \ensuremath{\mathcal{K}}} \ensuremath{\boldsymbol{S}}_{k, \ensuremath{\boldsymbol{q}}}
)^{\oplus\otimes} 
\circ D^2_{\ensuremath{\boldsymbol{q}}} \ensuremath{\psi}_t (\ensuremath{q}) \} + 
V_{\ensuremath{q}} \ensuremath{\psi}_t (\ensuremath{q}) 
. 
\end{align}
\end{subequations}
If $A$ is an operator acting on $E_{\ensuremath{\boldsymbol{q}}}$, then by $A^{\oplus \otimes}$ we
mean an operator defined to be the the direct sum over $\ensuremath{\nu} \in
\ensuremath{A_{q/K}}$ of the operators defined on the tensor product as $A$
acting on the \ensuremath{{\ensuremath{\boldsymbol{q}}}^{th}}\ component and the identity on the other
components. 
We can also write the evolution equations as 
\begin{subequations}
\begin{align} 
\ensuremath{\frac{d{\boldsymbol{Q}_t}}{d{t}}} =& \frac{\hbar}{\ensuremath{({\ensuremath{\psi}_t}, {\ensuremath{\psi}_t})}} 
\sum_{\ensuremath{\nu} \in \ensuremath{A_{q/K}}} \frac{1}{m_{\nu(\ensuremath{\boldsymbol{q}})}}  \textrm{Im}
\ensuremath{({\ensuremath{\psi}_t^{\ensuremath{\nu}}}, { \nabla_{\boldsymbol{Q}_t}
\ensuremath{\psi}_t^{\ensuremath{\nu}}})}   (Q_t) 
  \quad \text{for each $\boldsymbol{Q}_t \in Q_t$,} \\ 
\begin{split}
i \hbar \ensuremath{\frac{\partial {\ensuremath{\psi}_t^{\ensuremath{\nu}}}}{\partial {t}}}  (\ensuremath{q}) =&  - \frac{\hbar^2}{2} 
\sum\limits_{  \ensuremath{\boldsymbol{q}} \in 
  \ensuremath{q}} \{ \frac{1}{m_{\nu(\ensuremath{\boldsymbol{q}})}} \ensuremath{\Delta}_{\ensuremath{\boldsymbol{q}}} +
  \ensuremath{\boldsymbol{S}}_{\nu(\ensuremath{\boldsymbol{q}})} \circ (B_{\ensuremath{\boldsymbol{q}}}) \ensuremath{\psi}_t^{\ensuremath{\nu}} 
(\ensuremath{q}) \} 
+ 
(V_{\ensuremath{q}} \ensuremath{\psi}_t(\ensuremath{q}) )^{\ensuremath{\nu}} 
\end{split}
\end{align}
\end{subequations}
for each partition $ \ensuremath{\nu} \in \ensuremath{A_{q/K}}$.

\clearpage

\section{Appendix: the set-indexed tensor product}

\subsection{The set-indexed tensor product} \label{s:setten} \label{a:tensor}


Let $\ensuremath{T}$ be a finite set. For each element $\ensuremath{a} \in
\ensuremath{T}$, we associate 
a vector space $\ensuremath{W}_{\ensuremath{a}}$ to that element. We define the
\ensuremath{T}-indexed Cartesian product 
of these spaces to be 
\begin{equation*}
\ensuremath{\mathop{\times}_{\ensuremath{a} \in \ensuremath{T}}
\ensuremath{W}_{\ensuremath{a}} } \ensuremath{:=} \ensuremath{\{f | f: \ensuremath{T} \to \bigcup_{\ensuremath{a}
\in \ensuremath{W}}{\ensuremath{W}_{\ensuremath{a}} \text{ such that } f(\ensuremath{a}) \in
\ensuremath{W}_\ensuremath{a} \} }}.\footnote{We could also say that the Cartesian
product is the set of sections of 
the bundle over the discrete space $T$ whose fiber at $\ensuremath{a}$ is
$\ensuremath{W}_{\ensuremath{a}}$.  }
\end{equation*} 
Endowed with the usual vector space structure, this is the 
\ensuremath{T}-indexed direct
sum of the spaces.\footnote{
As an example, the tangent space of \ensuremath{{}^N\ensuremath{\mathbb{R}^{\ensuremath{3}}}}\ is the \ensuremath{q}-indexed direct sum of
the tangent 
spaces in physical space.
} 
We define the \ensuremath{T}-indexed tensor product \ensuremath{\bigotimes_{\ensuremath{a} \in \ensuremath{T}} \ensuremath{W}_{\ensuremath{a}} }\ to be the set of
multilinear
functionals defined on \ensuremath{\mathop{\times}_{\ensuremath{a} \in \ensuremath{T}}
\ensuremath{\ensuremath{W}^*}_{\ensuremath{a}} }, the \ensuremath{T}-indexed Cartesian product of the
dual spaces.    

If, as will be our case, the vector spaces $\ensuremath{W}_{\ensuremath{a}}$ are the
same vector space, say \ensuremath{W}, then simplifications
occur. Notationally, \ensuremath{{\ensuremath{W}}^{\ensuremath{T}}}\ denotes the Cartesian
product; \ensuremath{{\ensuremath{W}}^{{\otimes}{\ensuremath{T}}}}\  denotes the tensor product. As for the definition of
the Cartesian product, the
union can be replaced with the vector space itself. The
Cartesian product becomes
\begin{equation*}
\ensuremath{{\ensuremath{W}}^{\ensuremath{T}}} = \{f| f: \ensuremath{T} \to \ensuremath{W}\}.
\end{equation*}
From this point on, we will assume that we have only one vector
space. The more general case does not substantially change any
statements below, but it is irrelevant to the purpose of this paper. 

Let us consider a different approach to \ensuremath{{\ensuremath{W}}^{{\otimes}{\ensuremath{T}}}}. We will use the
standard equivalence class approach to tensors. We start by
introducing some notation. For $f \in \ensuremath{{\ensuremath{W}}^{\ensuremath{T}}}$, we write
\ensuremath{\mathop{\otimes}_{\ensuremath{a} \in \ensuremath{T}} \ensuremath{w}_{\ensuremath{a}} }, where $\ensuremath{w}_{\ensuremath{a}} \ensuremath{:=} f(\ensuremath{a})$. 
Given a decomposition of \ensuremath{T}\ into the subsets $\ensuremath{T}_{\alpha}$, $\alpha
\in A$, for some $A$, 
then we
certainly take as true 
 \[\mathop{\otimes}_{\alpha \in A} (\mathop{\otimes}_{\ensuremath{a} \in
\ensuremath{T}_{\alpha}} \ensuremath{w}_{\ensuremath{a}})
= \ensuremath{\mathop{\otimes}_{\ensuremath{a} \in \ensuremath{T}} \ensuremath{w}_{\ensuremath{a}} } = f.
\]
A special case of the above, and one with a special abuse of notation,
is when $A$ is a two element subset and one of the sets in the
decomposition has only one element, say $\ensuremath{a}'$. Then we would write 
\ensuremath{(\mathop{\otimes}_{ \substack{\ensuremath{a}
 \in \ensuremath{T} \\
 \ensuremath{a} \neq \ensuremath{a}'}} \ensuremath{w}_\ensuremath{a}) \otimes {{w_{\ensuremath{a}'}} }} which is equal to \ensuremath{\mathop{\otimes}_{\ensuremath{a} \in \ensuremath{T}} \ensuremath{w}_{\ensuremath{a}} }\ as well as
$w_{\ensuremath{a}'} \otimes (\mathop{\otimes}_{\substack{\ensuremath{a} \in
\ensuremath{T} \\ \ensuremath{a} \neq \ensuremath{a}'}} \ensuremath{w}_{\ensuremath{a}}) $.
We can then view the 
\ensuremath{T}-indexed tensor product as the equivalence classes of all formal
linear combinations of the elements in the set \ensuremath{{\ensuremath{W}}^{\ensuremath{T}}}\ under the
relations 
\begin{align*}
c(\ensuremath{\mathop{\otimes}_{\ensuremath{a} \in \ensuremath{T}} \ensuremath{w}_{\ensuremath{a}} }) &= \ensuremath{(\mathop{\otimes}_{ \substack{\ensuremath{a}
 \in \ensuremath{T} \\
 \ensuremath{a} \neq \ensuremath{a}'}} \ensuremath{w}_\ensuremath{a}) \otimes {{c\ensuremath{w}_{\ensuremath{a}'}} }} 
\\
\ensuremath{(\mathop{\otimes}_{ \substack{\ensuremath{a}
 \in \ensuremath{T} \\
 \ensuremath{a} \neq \ensuremath{a}'}} \ensuremath{w}_\ensuremath{a}) \otimes {{x_{\ensuremath{a}'}} }} + \ensuremath{(\mathop{\otimes}_{ \substack{\ensuremath{a}
 \in \ensuremath{T} \\
 \ensuremath{a} \neq \ensuremath{a}'}} \ensuremath{w}_\ensuremath{a}) \otimes {{y_{\ensuremath{a}'}} }} &= \ensuremath{(\mathop{\otimes}_{ \substack{\ensuremath{a}
 \in \ensuremath{T} \\
 \ensuremath{a} \neq \ensuremath{a}'}} \ensuremath{w}_\ensuremath{a}) \otimes {{(x
+y)_{\ensuremath{a}'}} }}.  
\end{align*}
These relations should hold for every $\ensuremath{a}' \in
\ensuremath{T}$, $c \in \ensuremath{\mathbb{C}}$ and
 $\ensuremath{w}_\ensuremath{a}$ $ \ensuremath{w}_{\ensuremath{a}'}$,  $x_{\ensuremath{a}'}$, $y_{\ensuremath{a}'} \in
\ensuremath{W}$.    
We should note the standard fact that
the multilinear functionals acting on \ensuremath{{\ensuremath{W}}^{\ensuremath{T}}}\ extend to
linear functionals acting on \ensuremath{{\ensuremath{W}}^{{\otimes}{\ensuremath{T}}}}; this is a third way of viewing \ensuremath{{\ensuremath{W}}^{{\otimes}{\ensuremath{T}}}}.
 
If \ensuremath{W}\ has an inner product, then there is a natural inner
product on \ensuremath{{\ensuremath{W}}^{{\otimes}{\ensuremath{T}}}}. Namely, for elements in \ensuremath{{\ensuremath{W}}^{\ensuremath{T}}}, we define the inner
product to be 
\begin{equation*}\ensuremath{({\ensuremath{\mathop{\otimes}_{\ensuremath{a} \in \ensuremath{T}} \ensuremath{w}_{\ensuremath{a}} }}, {\ensuremath{\mathop{\otimes}_{\ensuremath{a} \in
\ensuremath{T}} {{v}_{\ensuremath{a}} }}})} \ensuremath{:=}
\prod_{\ensuremath{a} 
\in \ensuremath{T}} \ensuremath{({\ensuremath{w}_\ensuremath{a}}, {v_{\ensuremath{a}}})}.
\end{equation*} 
We then extend this definition to \ensuremath{{\ensuremath{W}}^{{\otimes}{\ensuremath{T}}}}\ by using the appropriate
linearities.  

 If we have a linear operator $C$ acting on \ensuremath{W}, then
$C^{\ensuremath{a}}$ shall denote a linear operator acting on \ensuremath{{\ensuremath{W}}^{{\otimes}{\ensuremath{T}}}}\ such
that, when acting 
on product elements, it applies $C$ to the 
  $\ensuremath{a}^{th}$
factor while applying the identity to the other factors, i.e.\ 
\begin{equation*}
C^{\ensuremath{a}'} (\ensuremath{\mathop{\otimes}_{\ensuremath{a} \in \ensuremath{T}} \ensuremath{w}_{\ensuremath{a}} }) \ensuremath{:=} \ensuremath{(\mathop{\otimes}_{ \substack{\ensuremath{a}
 \in \ensuremath{T} \\
 \ensuremath{a} \neq \ensuremath{a}'}} \ensuremath{w}_\ensuremath{a}) \otimes {{C w_{\ensuremath{a}'}} }}
\end{equation*} 

We finish with a word about contractions. The set-indexed tensor product
is excellent for use with contractions. Indeed, given $w \in
\bigotimes_{a \in \alpha} E \otimes \bigotimes_{b \in \beta} E^{*}$,
then the contraction of $w$ over the index pair $(\tilde a, \tilde b)$
is denoted by $\ensuremath{\mathcal{C}}(\tilde a, \tilde b) w$ and is an element of
$\bigotimes_{a \in \{\alpha \setminus \{\tilde a\} \}} 
E \otimes \bigotimes_{b \in \{\beta \setminus \{\tilde b \}\} }
E^{*}$. This allows one to notationally keep track of which indices
were contracted. Using numbers as indices eliminates this
information. Furthermore, one can also see that a type of unit
analysis may be invoked. In other words, the expression $
\ensuremath{\mathcal{C}}(\tilde a, \tilde b) w + \ensuremath{\mathcal{C}}(a', {b'}) w$ would seem to make
no sense whereas 
$\ensuremath{\mathcal{C}}(\tilde a, \tilde b) \ensuremath{\mathcal{C}}(a',  b') w + 
\ensuremath{\mathcal{C}}(\tilde a,  b') \ensuremath{\mathcal{C}}( a', \tilde b) w$ would make sense. With
numbers, this would be hard to indicate. As such, the set-indexed tensor
product may be as
helpful in formulating meaningful expressions
 as dimensional analysis has historically been.

\clearpage

\section{Appendix: triviality of the Bose and Fermi bundles}

\subsection{On the 
question of the triviality of the Bose and Fermi
bundles} 
\label{a:triviality} \label{s:triv}


The notation of this section is that \ensuremath{\mathcal{B}(N, s)}\ is the Bose bundle
for spin-$s$ particles and \ensuremath{\ensuremath{\mathcal{F}}(N,s)}\ is the Fermi bundle. This is
the notation for $\ensuremath{d} =3$. For $\ensuremath{d} > 3$, we replace $s$ with the
dimension of the complex vector space. So $s$ is replaced with $2s
+1$. 

Here are the results of this section: 
\begin{enumerate} 
\item For $N =2$, $\ensuremath{d} =3$,  the triviality of the bundles are as follows:  

\begin{tabular}{|l|c|c|c|c|}
\hline
& $s=2l$  &  $s=2l+\frac{1}{2}$ &  $s=2l+1$ &  $s=2l+\frac{3}{2}$\\
\hline
$\ensuremath{\mathcal{B}(2,s)}$ & trivial & nontrivial & nontrivial & trivial \\
$\ensuremath{\ensuremath{\mathcal{F}}(2,s)} $ & nontrivial & nontrivial & trivial & trivial\\
\hline
\end{tabular}
for $l \in \mathbb{N}$. 
\item For $N \geq 2$, we have \ensuremath{\mathcal{B}(N, 2l+1)}\ is nontrivial and
\ensuremath{\ensuremath{\mathcal{F}}(N, 2l)}\ is nontrivial. We do not know the 
triviality or nontriviality of the other bundles.  
\item The configuration space of \ensuremath{{}^N\ensuremath{\mathbb{R}^{\ensuremath{d}}}}\ is orientable if and only if
\ensuremath{d}\ is even. 
\end{enumerate}


Throughout, $\ensuremath{d} \geq 3$. The nontriviality results apply for any
such \ensuremath{d}. The triviality results only apply for $N=2, \ensuremath{d}=3$ and for
the Bose line in any dimension. 


\begin{lemma}
A bundle is trivial iff it has a flat connection with trivial
holonomy. 
\end{lemma}

\begin{proof}
If it is trivial, say $M \times V$, then any basis of $V$ translates
into a global basis. Hence, the coordinates provide the way to do 
trivial transport. If a bundle has trivial holonomy, then
construct a basis by parallel transporting it. As parallel transport
along any closed path is, by assumption, the identity operation, the
constructed basis is well-defined and smooth. Note that having trivial
holonomy is the only way that a connection can  have a global parallel basis. 
\end{proof}


\begin{theorem}
If $L$ is a real or complex line bundle over $\ensuremath{\mathcal{Q}} \ensuremath{:=} \ensuremath{{}^N\ensuremath{\mathbb{R}^{\ensuremath{d}}}}$, then it is
nontrivial if there exists a flat connection $\nabla^L$ on $L$ 
such that its
holonomy is the alternating character of \ensuremath{\pi_1 ({\ensuremath{\mathcal{Q}}})}.
\end{theorem} 

\begin{proof}
The heart of the proof is the following idea. The difference between
two connections is a 1-form and for both of them to be flat, the
1-form must be closed. Furthermore, to have a nontrivial holonomy, the
1-form must be non-exact. Since all closed 1-forms on \ensuremath{\mathcal{Q}}\ are exact, we
are done. 

The details are the following.
We need to recall a few basic bundle facts. 
\begin{enumerate}
\item The difference of
two connections on the same bundle is a matrix-valued
1-form. For complex (real, respectively) line bundles, 
this means complex-valued
(real-valued) 1-forms. 
All we need to check is that it satisfies 
 \[ (\nabla^1 - \nabla^2 )(f \ensuremath{\psi}) = \nabla^1 f\ensuremath{\psi} -
\nabla^2 f\ensuremath{\psi} = f(\nabla^1 \ensuremath{\psi} -
\nabla^2 \ensuremath{\psi}),\]
i.e.\ it needs to be a tensor.
 This follows from the  Leibniz rule imposed on all connections;
 namely  
$\nabla f\ensuremath{\psi} = df \ensuremath{\psi} + f \nabla \ensuremath{\psi}$.    
\item Let $P_{\gamma}^2$ and $P_{\gamma}^1$ be the holonomy operators 
for $\nabla^1$, $\nabla^2$ along the path $\gamma$. By definition,
$\nabla^i_{\dot{\gamma}} P_{\gamma}^i = 0$. We define $C_{\gamma}$ to
satisfy  $P_{\gamma}^2 = C_{\gamma}
P_{\gamma}^1$. In particular, $C_{\gamma} = P_{\gamma}^2
(P_{\gamma}^1)^{-1}$.
If we further define $\omega \ensuremath{:=}
\nabla^2 - \nabla^1$ then we have 
\begin{equation} \label{se:cg}
\nabla^1_{\dot{\gamma}} C_{\gamma} = -\omega(\dot{\gamma}) C_{\gamma}.
\end{equation}
\item If a connection is flat, then  its parallel transport
operators depend only on the homotopy
class of the curve. 
\item For a line bundle, $\omega$ is a 1-form and $C_{\gamma}$ is
just a number. Then the solution to equation
\eqref{se:cg} is $C_{\gamma} = e^{-\int_{\gamma} \omega}$.
If we   
define $f(\gamma) = -\int_{\gamma} \omega$ and
if $\gamma$ and   $\sigma$ are in the same
homotopy class, then $f(\gamma)- f(\sigma)$  must be an integral
multiple of $2\pi i$   
in order that $C_{\gamma} = C_{\sigma}$. But the difference function
 is a continuous function
with a discrete set of values. Hence, it must be constant on any
connected component. In particular, the integrals depend only on the
homotopy class.    
\end{enumerate}

Assume that $L$ is trivial and let $\nabla^T$ be the implied trivial flat
connection on $L$.  Then $\omega$ defined by 
$\omega \ensuremath{:=} (\nabla^L - \nabla^T)$  is a 1-form on \ensuremath{\mathcal{Q}}. We have 
$P_{\gamma}^L = e^{\int_{\gamma} -\omega}P_{\gamma}^T $. 
By
assumption, if $\ensuremath{\sigma}$ is a path in \ensuremath{\mathcal{Q}}\ exchanging two
particles, then ${\int_{\ensuremath{\sigma}} -\omega} \equiv i \pi (\text{mod }
2\pi),$ i.e. it is non-zero.  But 
\[
0 
= \int_{\text{Id}} -\omega 
= \int_{\ensuremath{\sigma}\ensuremath{\sigma}} -\omega 
= 2 \int_{\ensuremath{\sigma}} -\omega. 
\]
Thus, we have that the integral must be zero which contradicts
nontrivial holonomy.
\end{proof}
 
\begin{cor}
The Fermi line on \ensuremath{{}^N\ensuremath{\mathbb{R}^{\ensuremath{d}}}}\ is nontrivial. 
\end{cor}

\begin{proof}
The Fermi line's holonomy is the alternating character. Thus it is
a nontrivial line bundle.
\end{proof}

\begin{cor}
The configuration space of \ensuremath{{}^N\ensuremath{\mathbb{R}^{\ensuremath{d}}}}\ is orientable if and only if \ensuremath{d}\ is
even. 
\end{cor}

\begin{proof}
A space is orientable if and only if the determinant line of the
tangent space is trivial. That is to say, we need to compute
$\Lambda^{\ensuremath{d} N} (T\ensuremath{\mathcal{Q}})$. Let $\{e_i\}_{i=1}^{\ensuremath{d}}$ be the standard
basis in \ensuremath{\mathbb{R}^{\ensuremath{d}}}. Then a basis element of $\Lambda^{\ensuremath{d} N}
(T\ensuremath{\mathcal{Q}})$ is $\mathop{\wedge}\limits_{i=1}^{\ensuremath{d}} (\mathop{\wedge}\limits_{\ensuremath{\boldsymbol{q}} \in
\ensuremath{q}} e_{i,\ensuremath{\boldsymbol{q}}}) $; this is defined up to a sign. In order to
define a global basis, we must be able to choose the sign
consistently. Rather than do that, we shall compute the holonomy of
this bundle:
\[
 \mathop{\wedge}_{i=1}^{\ensuremath{d}} (\mathop{\wedge}_{\ensuremath{\boldsymbol{q}} \in \ensuremath{q}} e_{i,\ensuremath{\boldsymbol{q}}}) 
\overset{\ensuremath{\Gamma_{\ensuremath{\sigma}}}}{\mapsto}  
\mathop{\wedge}_{i=1}^{\ensuremath{d}} (\mathop{\wedge}_{\ensuremath{\boldsymbol{q}} \in \ensuremath{q}} e_{i,\ensuremath{\sigma}\ensuremath{\boldsymbol{q}}})  
= (-1)^{\ensuremath{d} |\ensuremath{\sigma}|}\mathop{\wedge}_{i=1}^{\ensuremath{d}} (\mathop{\wedge}_{\ensuremath{\boldsymbol{q}} \in \ensuremath{q}} e_{i,\ensuremath{\boldsymbol{q}}}).
\]
That is to say, the induced flat connection has trivial holonomy if
and only if \ensuremath{d}\ is even. By the above results, we therefore have that 
$\Lambda^{\ensuremath{d} N} (T\ensuremath{\mathcal{Q}})$ is trivial if and only if \ensuremath{d}\ is
even. 
\end{proof}

\begin{lemma}
The following relations hold: $\ensuremath{\mathfrak{b}} \otimes \ensuremath{\mathfrak{b}} =
\ensuremath{\mathfrak{b}}$,  $\ensuremath{\mathfrak{f}} \otimes \ensuremath{\mathfrak{f}} =
\ensuremath{\mathfrak{b}}$, $\ensuremath{\mathfrak{f}} \otimes \ensuremath{\mathfrak{b}} =
\ensuremath{\mathfrak{f}}$
\end{lemma}

\begin{proof}
It immediately follows from the fact that the holonomy operators of a
tensor product are the tensor products of the original holonomies.
In other words, the characters multiply under tensoring of the bundles. 
\end{proof}

The following statement and argument is for $\ensuremath{d} =3$. 

\begin{lemma}
Let $F$ be the Whitney sum (direct sum) 
of the 2 particle Fermi line with itself. 
Then $F = \ensuremath{\mathfrak{f}} \oplus \ensuremath{\mathfrak{f}} \ensuremath{\simeq} {}^2\mathbb{R}^3
\times \ensuremath{\mathbb{C}}^2$. 
\end{lemma} 


\begin{proof}
We start by coordinatizing the covering space of $\ensuremath{\mathcal{Q}}
\ensuremath{:=}{}^2\mathbb{R}^3 $: 
\[
 \ensuremath{\widehat{\ensuremath{\mathcal{Q}}}} \ensuremath{:=} \ensuremath{\mathbb{R}^{2,3}}_{\neq} = \ensuremath{\mathbb{R}^{3}} \times \ensuremath{\mathbb{R}^{+}} \times S^2.
\]
 That is to say, we coordinatize the particle system  by the
center of mass, the distance between the two particles, and the
direction from particle 1 to particle 2. To create the base space, we
identify antipodal points on $S^2$ which corresponds to not knowing which
particle is particle 1. Our plan is to create two non-zero sections of the
trivial bundle  $\ensuremath{\widehat{\ensuremath{\mathcal{Q}}}} \times \ensuremath{\mathbb{C}}^2$ that satisfy
the alternating periodicity condition. We first do this 
on the sphere and then extend it to the whole space by
constancy in the other coordinates. 

Although we shall write it down explicitly, the map from the sphere
into sections of $\ensuremath{\mathbb{C}}^2$ will be the procedure of taking the
columns of the unitary matrix which represents, in the
spin-$\frac{1}{2}$ representation, a rotation about the line in the
sphere, by an angle $\pi$. The point is that the representation is a
double cover of the rotation group and the two differ by a minus
sign. Since rotation about $\pi$ is the same rotation element for a
point on the sphere and its antipode, we expect that this is what we
want. Indeed, let the usual representation with the Pauli spin
matrices, as generators of the representation, be given. 
Then a rotation about an angle $\alpha$ about the line specified by
the point in the sphere, written as a $(x, y, z)$ vector with the
sphere  parametrized 
 by the angles $0 \leq \theta \leq 2\pi$ and $0
\leq \phi \leq \pi$, is
$\boldsymbol{p}= (\sin\phi\cos\theta, \sin\phi\sin\theta, \cos\phi)$ 
and it  maps to the unitary matrix, expressed in the basis of $\sigma_{z}$
eigenstates,  
\[ 
\begin{split}
e^{i \frac{\alpha}{2} \boldsymbol{\sigma} \cdot \boldsymbol{p}} 
&= \cos\frac{\alpha}{2} + 
i\boldsymbol{\sigma} \cdot \boldsymbol{p}\sin\frac{\alpha}{2} \\
& =
\cos\frac{\alpha}{2} + i\sin\frac{\alpha}{2} \begin{pmatrix}
\cos\phi & \sin\phi\cos\theta - i\sin\phi\sin\theta \\ 
\sin\phi\cos\theta - i\sin\phi\sin\theta &\cos\phi
\end{pmatrix}.
\end{split} 
\]
Thus, if we use the angle $\alpha = \pi$ and take the columns as the
sections, we have, on $S^2\times \ensuremath{\mathbb{C}}^2$, 
\[
\ensuremath{\psi}_1 (\theta, \phi)= 
\begin{pmatrix}
\cos\phi\\
e^{i\theta}\sin\phi
\end{pmatrix}
\quad
\ensuremath{\psi}_2 (\theta, \phi) = 
\begin{pmatrix} 
e^{-i\theta} \sin\phi\\
-\cos\phi
\end{pmatrix}
\]  
Being the columns of a unitary matrix, or by direct computation, 
the norm of $\ensuremath{\psi}_1$ and $\ensuremath{\psi}_2$ are both 1,
they are orthogonal vectors, they are smooth 
sections,
and they satisfy 
\[
\ensuremath{\psi}_i(\theta + \pi, \pi - \phi) = - \ensuremath{\psi}_i(\theta, \phi).
\] 
More to the point, if we define a new bundle by 
$(p, v) \sim (-p, -v)$ where $p$ and $-p$ are antipodal points, 
then $[\ensuremath{\psi}_1]$ and $[\ensuremath{\psi}_2]$ form
orthonormal sections of the new bundle.\footnote{There is another
tempting map which is to take the eigenspinors of the spin angular
momentum operator in the direction represented by $p$. This fails
since under antipodal exchange, the vectors are exchanged. We wanted
to only change the sign.} 

To go back to the original problem, we define $\Psi_i (x, r, [p]) 
\ensuremath{:=} [\ensuremath{\psi}_i]([p])$. These are sections of  $F$. Since
they form a basis at every point, $F$ is trivial. 
\end{proof}

\begin{lemma}
Let $C$ be the Whitney sum of the Fermi line for $N$ particles with
the  trivial
bundle $\ensuremath{{}^N\ensuremath{\mathbb{R}^{\ensuremath{d}}}} \times \ensuremath{\mathbb{C}}^l$. 
Then  $C$ is not a trivial bundle, i.e. 
$ C  = \ensuremath{\mathfrak{f}} \oplus \ensuremath{\mathbb{C}}^l \ensuremath{\not\simeq} \ensuremath{{}^N\ensuremath{\mathbb{R}^{\ensuremath{d}}}} \times
\ensuremath{\mathbb{C}}^{l+1}$.
\end{lemma} 

\begin{proof}
This follows from the fact that the determinant line of a tensor
product is the tensor product of the determinant lines, i.e 
$\Lambda^{mn}(E \otimes F) \ensuremath{\cong} \Lambda^{n}(E)\otimes
\Lambda^{m}(F)$ where $E$ has dimension $n$ and $F$ has dimension
$m$.\footnote{This follows from the general identity 
\[ 
\Lambda^k(E\otimes F) \ensuremath{\cong} \sum_{i=0}^k \Lambda^i(E) \otimes
\Lambda^{k-i}(F)
\]    
where $\Lambda^0(\cdot)$ is defined to be the trivial line bundle.} 
Since the determinant line of a trivial bundle is \ensuremath{\mathfrak{b}}, we have that
the determinant line of $C \ensuremath{\cong} \ensuremath{\mathfrak{f}} \otimes \ensuremath{\mathfrak{b}}$. 
That is to say, $\Lambda^{l+1}(C) \ensuremath{\cong}
\ensuremath{\mathfrak{f}}$. Hence it is nontrivial as the determinant line is
nontrivial and any trivial bundle has a trivial determinant line.  
\end{proof}

\begin{theorem}
For $N =2$,  the triviality of the bundles are as follows:  

\begin{tabular}{|l|c|c|c|c|}
\hline
& $s=2l$  &  $s=2l+\frac{1}{2}$ &  $s=2l+1$ &  $s=2l+\frac{3}{2}$\\
\hline
$\ensuremath{\mathcal{B}(2,s)}$ & trivial & nontrivial & nontrivial & trivial \\
$\ensuremath{\ensuremath{\mathcal{F}}(2,s)} $ & nontrivial & nontrivial & trivial & trivial\\
\hline
\end{tabular}
for $l \in \mathbb{N}$. 
\end{theorem}

\begin{proof}
The idea of this proof is to decompose each bundle into a sum of Bose
and Fermi lines. We then pair the Fermi lines to make trivial
bundles. We are then either left with a completely trivial bundle or
the sum of a trivial bundle with a single Fermi line. 

Let us start with the Bose bundles. Since we have only two particles,
the bundle is the sum of a symmetric part and an antisymmetric
part. The symmetric subbundle is trivial (a sum of Bose lines), 
while the antisymmetric part
is a sum of Fermi lines. Thus, the Bose bundle is trivial if and only
if the number of Fermi lines is even, i.e. if the dimension of the 
antisymmetric part is even. 
As for the
Fermi bundle, its triviality depends upon the parity of the dimension of the
symmetric part of the Bose bundle 
since upon tensoring with a Fermi line, the Fermi lines
become Bose lines and the Bose lines become Fermi lines. 

More specifically, the Bose bundle \ensuremath{\mathcal{B}(N, s)}\
 may be written as, using the basis
$\ensuremath{w}^i$ of $\ensuremath{\mathbb{C}}^{2s+1}$,  
\[\ensuremath{\mathcal{B}(N, s)} 
= 
\bigoplus_{i=1}^{2s+1} 
sp\{\ensuremath{w}_{\ensuremath{q}}^i \otimes \ensuremath{w}_{\ensuremath{q}'}^i \} 
\oplus 
\bigoplus_{1 \leq i < j \leq 2s+1} 
(sp\{
\ensuremath{w}_{\ensuremath{q}}^i \otimes \ensuremath{w}_{\ensuremath{q}'}^j
+ 
\ensuremath{w}_{\ensuremath{q}}^j \otimes \ensuremath{w}_{\ensuremath{q}'}^i
\}
\oplus 
sp\{
\ensuremath{w}_{\ensuremath{q}}^i \otimes \ensuremath{w}_{\ensuremath{q}'}^j
-
\ensuremath{w}_{\ensuremath{q}}^j  \otimes \ensuremath{w}_{\ensuremath{q}'}^i
\}) 
\]
where $sp\{v\}$ means the linear span of the set of vectors; in this
case these are all 1-dimensional subspaces.
Thus it is the sum of $2s+1 + \binom{2s+1}{2}$ Bose lines and
$\binom{2s+1}{2}$ Fermi lines. Since $\binom{2s+1}{2}$ is odd exactly
when $s = 2l + \frac{1}{2}$ and $s= 2l+1$, we have the Bose bundle
results.
 For the Fermi bundle
$\ensuremath{\ensuremath{\mathcal{F}}(N,s)} (= \ensuremath{\mathfrak{f}} \otimes \ensuremath{\mathcal{B}(n,s)})$ we have that it is a sum
of 
$\binom{2s+1}{2}$  Bose lines
and $2s+1 + \binom{2s+1}{2}$   Fermi lines as 
tensoring distributes
over sums. Since $2s+1 + \binom{2s+1}{2}$ is odd exactly when $s= 2l$
and  $s=2l+\frac{1}{2}$, we have our results. 
\end{proof}

 
\begin{lemma}
The determinant line $\Lambda^n(E)$ of a flat bundle $E$ 
has an induced connection whose
holonomy representation is  multiplication by the determinant of
the 
corresponding holonomy operator on $E$. 
\end{lemma} 

\begin{proof}
The claim follows immediately
from the standard fact  that, for any matrix $A:V^n \to V^n$, we have
that $\det A$ satisfies 
\[ 
\mathop{\wedge}_{i=1}^n A e_i = \det A (\mathop{\wedge}_{i=1}^n e_i)
\]
where each $e_i \in V$. Thus, if $\{e_i\}_{i=1}^n$
forms a basis, then $\det A$ is determined by the above equation. 


Thus, the determinant line's induced holonomy is
given by the determinant map of the holonomy representation.  
\end{proof} 

\begin{theorem}
 For $N \geq 2$ and $\ensuremath{d} \geq 3$, we have that $(\ensuremath{\mathbb{C}}^{4k +2})^{\otimes
 \ensuremath{\mathcal{Q}}}$ and $\ensuremath{\mathfrak{f}} \otimes (\ensuremath{\mathbb{C}}^{4k +1})^{\otimes
 \ensuremath{\mathcal{Q}}}$ are nontrivial bundles.
 

\end{theorem}  

\begin{proof}
We will establish that their determinant lines are nontrivial. In
fact, we will show that the determinant of the parallel transport
operator which exchanges two particles is $-1$. By the above lemma and
the first theorem, 
that is sufficient to establish our result. Fix the two positions to
be exchanged. Let $l$ be the dimension of the factor in the tensor
product. Ignoring the Fermi line factor, 
symmetrize and antisymmetrize on the two factors that
are being exchanged. Count how many of each kind there are. The
total number of symmetric elements is $(l +
\binom{l}{2})*(l)^{N-2}$ while the total number of antisymmetric
elements is $\binom{l}{2}*(l)^{N-2}$. Thus, the bundle without the
Fermi line involved is
nontrivial if $l$ is odd and $\binom{l}{2}$ is
odd. This occurs when $l = 4k+2$; for $N=2$, we also have $l = 4k+3$. 
Tensoring with the Fermi line 
means the resultant bundle is nontrivial
if $l$ is odd and $(l +
\binom{l}{2})$  is odd. This occurs when $l = 4k+2$; for
$N=2$, we also have $l= 4k+1$. 
We are done.
\end{proof}

\clearpage
\bibliographystyle{abbrv}
\bibliography{references}
\end{document}
